% Propietario conservado: /Users/ruben/Documents/New project/output/EXTREMA_MEDIA_RAZON_RECIPROCIDAD_ES_EN_20260918/ARTICULO/ES/source/sections/generacion.tex
\section{Génération coinductive de \texorpdfstring{\(\pi,\varphi,e\)}{pi, phi, e} à profondeur arbitraire}
\label{x_reciprocidad:sec:generacion}

L'émission de six blocs est une étape finie de la construction, et non un développement décimal complet. Sa fonction consiste à produire des régions avec leur multiplicité, leur orientation et leurs règles de transport. Le prolongement doit conserver ces données et déterminer une famille compatible de préfixes de longueur arbitraire. Le mot \emph{génération} désigne cette opération ; \emph{généalogie} désigne la reconstruction de ses dépendances. Les deux notions interviennent, mais remplissent des fonctions différentes.

\subsection{Orbites, régions et caractères de clôture, de propagation et d'auto-échelle}

Sur un mot \(w=(w_1,\ldots,w_6)\), nous définissons
\[
 r(w)=(w_3,w_1,w_2,w_5,w_6,w_4),\qquad
 s(w)=(w_4,w_5,w_6,w_1,w_2,w_3).
\]
On a \(r^3=s^2=I\) et \(srs=r^{-1}\) ; ces permutations réalisent donc le groupe diédral \(D_3\). L'action conserve le catalogue et les multiplicités et commute avec \(\Psi\). Le quotient des \(243\) mots visibles possède \(43\) orbites : trente-huit de cardinal six et cinq de cardinal trois.

La classe de clôture se reconnaît par un prédicat de fibre : ses six orientations ont cinq émissions par mot, une multiplicité totale \(1008\) et une distribution \(432+4\cdot144\). Il existe une seule orbite possédant ces propriétés dans le catalogue. Pour préciser les deux autres sélections, nous écrivons \(f(w)=\#\Psi^{-1}(w)\), \(m(w)=\sum_{U\in\Psi^{-1}(w)}\operatorname{mult}(U)\) et \(\nu(w)=(n_0,n_1,n_2)\). La classe diagonale additive d'une émission est \(d(U)=[i+j]_9\) dans les indices de zéro à huit ; sa signature additive conserve cette classe et l'orientation \(ES\) ou \(NO\). Soit \(\mathcal D(\mathcal O)\) l'ensemble de ces classes pour toutes les émissions sur une orbite.

Parmi les orbites de taille six, le prédicat conjoint
\[
 f(w)=1,\qquad m(w)=144\quad(w\in\mathcal O),\qquad
 \mathcal D(\mathcal O)=\{0,3,6\}
\]
laisse exactement six orbites. Dans ce secteur, le recensement \(\nu=(2,3,1)\), égal à celui de clôture, sélectionne uniquement la propagation ; le recensement \(\nu=(2,2,2)\) sélectionne uniquement l'auto-échelle. Le support diagonal ne peut pas être omis : sans lui, il y a quatre orbites à un seul point équilibrées de multiplicité \(144\), et non une seule. Enfin, l'existence d'une émission avec \(d=6\) et une orientation \(ES\) sélectionne un représentant de chaque orbite et conduit à
\begin{equation}
 w^{\rm cl}=010211,\qquad w^{\rm pr}=201101,\qquad w^{\rm au}=121200.
 \label{x_reciprocidad:eq:palabras-regionales}
\end{equation}
Ces symboles désignent d'abord des régions du catalogue. L'identification scalaire sera démontrée au moyen de leurs lecteurs ; ils ne sont pas choisis parce qu'ils coïncident avec des chiffres d'une constante connue.

Le mot de clôture \(010211\) possède les cinq préimages suivantes :
\begin{center}\small
\begin{tabular}{@{}lll r@{}}\toprule
Émission de six blocs&Firma multiplicativa&Rôle&Multiplicité\\\midrule
\(501,614,498,169,272,272\)&\(555555\)&transversal&432\\
\(810,923,870,169,272,272\)&\(339555\)&orientation positive&144\\
\(810,983,810,169,272,272\)&\(393555\)&orientation positive&144\\
\(870,923,810,169,272,272\)&\(933555\)&orientation positive&144\\
\(870,923,810,418,674,521\)&\(933717\)&retour orienté&144\\\bottomrule
\end{tabular}
\end{center}
La multiplicité \(432\) et la signature constante \(555555\) caractérisent la région transversale. Elles n'identifient pas la position centrale \((5,5)\) d'APP. Confondre une signature de six fenêtres avec une position du support éliminerait la distinction entre région numérique et structure électronique.

La microfibre possède en outre une morphologie bipartite précise. Chaque émission est un produit d'une classe de curseur additif et d'une classe multiplicative. Les trois régions positives partagent la même classe additive de dix-huit conditions ; ensemble, elles contiennent trois classes multiplicatives de huit conditions chacune. La région transversale possède un produit \(18\times24\) ; les trois positives réunies, un autre \(18\times24\) ; et le retour, un produit \(18\times8\). La décomposition par émissions comporte cinq parties, tandis que la décomposition en composantes connexes du graphe biparti en comporte trois. Toutes deux conservent les mêmes \(1008\) germes sans réduire la forme à un cardinal.

% Procedencia: inventario REV05, SR-015--027; ampliacion 04_tpk_emisor,
% firmas y factorizacion; X_RECIPROCIDAD/sections__generacion.tex:16--43.
% Ampliacion de la prueba de metadatos y de las tres componentes.
\subsection{Métadonnées orientées et décomposition de la microfibre régionale}
\label{sec:microfibra-bipartita-explicita}

La signature additive conserve la classe diagonale d'origine et l'orientation
de son déplacement. Pour un curseur initial $(i,j;d)$, avec des indices entre
zéro et huit, nous définissons
\[
 c=[i+j]_9,\qquad
 h(d)=\begin{cases}ES,&d\in\{E,S\},\\NO,&d\in\{N,O\}.
 \end{cases}
\]
Dans les activations additives consécutives, la classe $c$ augmente d'une unité
pour $ES$ et diminue d'une unité pour $NO$. Après la neuvième activation, la
direction est inversée. La suite des échantillons est donc déterminée
par $(c,h)$, avec l'inversion après le neuvième échantillon. L'échantillon
dans la classe $c'$ est $\rho_9(c'+2)$ ; chaque signature décimale additionne trois échantillons
consécutifs et réduit le total modulo dix.

\begin{proposition}[Récupération des métadonnées additives]
Le lecteur $(c,h)\mapsto f_+(c,h)$ est bijectif sur les dix-huit
signatures additives de l'émetteur. Chaque préimage dans l'espace des curseurs
comprend dix-huit conditions initiales. Les deux coordonnées sont
récupérées au moyen du tableau suivant.
\end{proposition}

\begin{center}
\begin{tabular}{c|cc}
\toprule
$c$&$f_+(c,ES)$&$f_+(c,NO)$\\
\midrule
0&\texttt{988212}&\texttt{212988}\\
1&\texttt{212645}&\texttt{645212}\\
2&\texttt{546988}&\texttt{988546}\\
3&\texttt{889221}&\texttt{221889}\\
4&\texttt{122564}&\texttt{564122}\\
5&\texttt{465898}&\texttt{898465}\\
6&\texttt{898122}&\texttt{122898}\\
7&\texttt{221456}&\texttt{456221}\\
8&\texttt{654889}&\texttt{889654}\\
\bottomrule
\end{tabular}
\end{center}

\begin{proof}
Pour chaque entrée, on applique la récurrence de classes décrite avant
l'énoncé pendant dix-huit activations ; leurs groupes de trois donnent les six
chiffres imprimés. Par exemple, depuis $(c,h)=(4,NO)$, les neuf premières
classes sont $4,3,2,1,0,8,7,6,5$ et les échantillons sont
$6,5,4,3,2,1,9,8,7$. Les trois totaux sont $15,6,24$, qui publient
$5,6,4$. L'inversion laisse la seconde moitié $1,2,2$, ce qui donne
\texttt{564122}. Les dix-huit chaînes du tableau sont distinctes,
de sorte que le lecteur est injectif et énumère toute son image. Pour chaque
$c$, il existe neuf couples $(i,j)$, un pour chaque $i$, et pour chaque $h$, il existe
deux directions. Cela donne $9\cdot2=18$ conditions par signature et
$18\cdot18=324$ conditions au total.
\end{proof}

Soient $A_s=f_+^{-1}(s)$ et $B_e=f_\times^{-1}(e)$ les classes de curseurs
des deux feuillets. La loi décimale récupère les deux signatures de chaque émission
$U$ : si $U_j=100s_j+10t_j+u_j$, alors
$s_j=\lfloor U_j/100\rfloor$ et $e_j=[t_j-s_j]_{10}$.
Par conséquent, la préimage de $U$ est exactement $A_s\times B_e$.
Les conditions initiales demeurent dans ce produit ; leurs cardinaux
constituent une lecture supplémentaire de la même préimage.

\begin{definition}[Incidence bipartite de la microfibre]
Pour le mot orienté $w=\texttt{010211}$, on forme le graphe dont les
sommets d'une classe sont les signatures additives qui apparaissent au-dessus de $w$ et
dont les sommets de l'autre classe sont les signatures multiplicatives. Une
arête $(s,e)$ représente l'émission correspondante et conserve comme
préimage étiquetée le produit $A_s\times B_e$.
\end{definition}

La récupération des signatures des cinq émissions donne le tableau complet
des arêtes :
\begin{center}
\begin{tabular}{llrrl}
\toprule
$s$&$e$&$|A_s|$&$|B_e|$&strate\\
\midrule
\texttt{564122}&\texttt{555555}&18&24&transversal\\
\texttt{898122}&\texttt{339555}&18&8&positif\\
\texttt{898122}&\texttt{393555}&18&8&positif\\
\texttt{898122}&\texttt{933555}&18&8&positif\\
\texttt{898465}&\texttt{933717}&18&8&retour\\
\bottomrule
\end{tabular}
\end{center}

\begin{theorem}[Décomposition exacte en trois composantes]
L'incidence bipartite possède trois composantes connexes, isomorphes à
$K_{1,1}$, $K_{1,3}$ et $K_{1,1}$. Leur préimage conjointe est l'union
disjointe
\begin{align*}
 &A_{\texttt{564122}}\times B_{\texttt{555555}},\\
 &A_{\texttt{898122}}\times
  (B_{\texttt{339555}}\sqcup B_{\texttt{393555}}
     \sqcup B_{\texttt{933555}}),\\
 &A_{\texttt{898465}}\times B_{\texttt{933717}}.
\end{align*}
Les cardinaux respectifs sont $432,432,144$ et leur somme vaut $1008$.
La décomposition par émissions comprend cinq produits et la
décomposition par composantes comprend trois ensembles.
\end{theorem}
\begin{proof}
Les trois signatures additives imprimées sont distinctes. Les cinq signatures
multiplicatives le sont aussi. Seules les trois arêtes positives partagent
un sommet, celui associé à \texttt{898122} ; elles forment une étoile connexe
à trois arêtes. Les deux autres arêtes n'ont aucun sommet commun avec
elle ni entre elles, et chacune forme une composante. Les préimages de
signatures différentes par une fonction sont disjointes. La factorisation de
l'émetteur et la distributivité du produit cartésien sur cette union
disjointe démontrent les trois expressions. Leurs cardinaux sont
$18\cdot24$, $18(8+8+8)$ et $18\cdot8$. Les cinq arêtes restent
individualisées au sein de ces composantes, avec leur multiplicité et leur
émission ordonnée.
\end{proof}

L'application de récupération des métadonnées donne $(c,h)=(4,NO)$ pour la région transversale,
$(6,ES)$ pour chacune des trois positives et $(5,NO)$ pour le retour.
Ainsi, l'orientation, la classe diagonale, l'incidence et la multiplicité se lisent dans
les mêmes signatures produites par les curseurs. La distinction entre cinq
émissions et trois composantes conserve la structure complète du recensement,
y compris la coïncidence du cardinal $432$ entre deux composantes distinctes.

\subsection{Transport fini et frontière de prolongement}

Les endomorphismes \(L_0,L_1\) de l'espace des mots ternaires agissent avant la réalisation archimédienne. Leur détermination commence par les transitions de l'état arithmétique, et non par la prescription de deux matrices. Si \(U_t\) est l'étape composée de sélection, de transport et de mise à jour, le bloc observé satisfait
\[
 b_j^\chi=\Pi_6(x_j^\chi),\qquad
 x_{j+1}^\chi=U_{9j+9}\cdots U_{9j+1}(x_j^\chi),
 \qquad \chi\in\{\mathrm{cl},\mathrm{pr},\mathrm{au}\}.
\]
On forme six transitions pour chaque régime : deux consécutives dans chacune des trois régions. Les couples de matrices d'origines et de destinations obtenus dans le registre sont
\[
\begin{array}{c|c|c|c}
X_0&Y_0&X_1&Y_1\\\hline
010211&012222&010211&002111\\
012222&010211&002111&110221\\
201101&121221&102011&012222\\
121221&102011&012222&102011\\
121200&112202&121020&010210\\
112202&121020&010210&010200
\end{array}
\]
Chaque chaîne de six chiffres représente une ligne dans \(\FF_3^6\). Puisque \(\det X_0=1\) et \(\det X_1=-1\) dans \(\FF_3\), les équations \(X_aL_a=Y_a\) déterminent univoquement \(L_a=X_a^{-1}Y_a\). Dans cette carte de vecteurs-lignes, on obtient
\[
 L_0=\begin{pmatrix}
 2&2&2&1&2&1\\2&1&2&2&1&1\\1&0&0&1&0&2\\
 0&0&1&1&1&0\\2&2&2&0&2&1\\2&1&2&1&0&0
 \end{pmatrix},\quad
 L_1=\begin{pmatrix}
 2&1&2&0&2&2\\1&2&1&1&2&0\\1&2&1&1&1&2\\
 0&1&0&0&2&1\\2&0&2&1&0&1\\0&2&2&2&1&1
 \end{pmatrix}\quad(\bmod3).
\]
Le calendrier \((L_0,L_0,L_1,L_1)\) produit quatre nouveaux blocs de six composantes. Leurs concaténations forment
\[
 w_6\longrightarrow w_{12}\longrightarrow w_{18}
 \longrightarrow w_{24}\longrightarrow w_{30}.
\]
Les sous-indices de \(w\) indiquent la longueur accumulée. La frontière \(R_{36}\) conserve la structure terminale et ouvre la relation de prolongement ; elle n'est pas renommée comme un dernier bloc périodique qui pourrait se répéter indéfiniment.

L'identité \(L=X^{-1}Y\) démontre l'unicité une fois les transitions fixées ; elle ne remplace pas la production de ces transitions par \(U_t\). La provenance de leur registre et l'examen des seize calendriers binaires sont documentés dans le supplément technique. Le calendrier indiqué est le seul de ces seize calendriers qui reproduise conjointement les trois suites de cinq blocs.

\subsection{Cinq régions de clôture et compensation angulaire}

La région de clôture contient une composante transversale et quatre accompagnatrices. Le lecteur symétrique agit dans l'espace de leurs cinq positions au moyen de
\[
 P_5=\frac15\mathbf1\mathbf1^{\mathsf T},\qquad
 \mathbf1=(1,1,1,1,1)^{\mathsf T}.
\]
L'invariance \(P_5\lambda=\lambda\) et la conservation \(\mathbf1^{\mathsf T}\lambda=1\) imposent \(\lambda=\mathbf1/5\). Cette normalisation répartit l'unité entre les positions du lecteur ; elle ne remplace pas les multiplicités de germes \(432,144,144,144,144\) par des fréquences uniformes. La carte elliptique du lecteur réalise chaque coordonnée au moyen de \(1+\lambda_jJ\) ; sa pente est \(\lambda_j\), de sorte que la pente primitive est \(1/5\). La règle qui passe d'une coordonnée normalisée à une pente est ainsi explicitée, au lieu d'identifier les deux notions sans application. Le lecteur compose les quatre positions accompagnatrices par la loi
\begin{equation}
 a\oplus b=\frac{a+b}{1-ab}.
 \label{x_reciprocidad:eq:suma-pendientes}
\end{equation}
Cette loi s'obtient en multipliant \((1+aJ)(1+bJ)\) dans le régime \(J^2=-I\) et en divisant par sa composante scalaire. La composition des pentes est donc une réalisation de la structure quadratique de TRIT, et non une somme ordinaire de quatre fractions.

\begin{proposition}[Compensateur de la clôture]
La composition de quatre pentes \(1/5\) est \(120/119\). La condition de retour à une pente un détermine un unique compensateur positif de la forme \(1/q\), avec \(q>1\), et donne \(q=239\).
\end{proposition}
\begin{proof}
La première duplication donne \((1/5)\oplus(1/5)=5/12\) ; la seconde, \((5/12)\oplus(5/12)=120/119\). La condition
\[
 \frac{120/119-1/q}{1+(120/119)/q}=1
\]
équivaut à \((120-119)q=120+119\), d'où \(q=239\).
\end{proof}

Le caractère de clôture résultant a pour réalisation
\begin{equation}
 \Pi_{\rm cl}=4\left(4\arctan\frac15-\arctan\frac1{239}\right).
 \label{x_reciprocidad:eq:lector-pi}
\end{equation}
L'identité angulaire de Machin reconnaît \(\Pi_{\rm cl}=\pi\). Le sens de la construction est précis : la pente et le nombre de compositions sont déclarés par le lecteur de la pentafibre ; l'équation de compensation impose \(239\). Une fois ce caractère produit, son évaluation par des séries alternées ne sélectionne pas à nouveau l'orbite des germes.

Pour \(0<x<1\), les sommes alternées
\[
 A_n(x)=\sum_{j=0}^n\frac{(-1)^jx^{2j+1}}{2j+1}
\]
encadrent \(\arctan x\) entre deux troncatures consécutives, dont la différence est \(x^{2n+3}/(2n+3)\). Appliquées aux deux pentes de \eqref{x_reciprocidad:eq:lector-pi}, elles fournissent des extrémités rationnelles d'erreur explicite. La valeur numérique est ainsi calculée comme conséquence du caractère exact, sans introduire de préfixe cible.

\subsection{Lecteur de propagation et normalisation de l'unité}

La région de propagation se réalise au moyen d'une famille formelle \(P_t\in\mathcal A[[t]]\), où \(\mathcal A\) est une algèbre commutative unitaire sur \(\QQ\). Sa composition satisfait \(P_{t+s}=P_tP_s\), avec pour unité \(P_0=1\), et le transport élémentaire orienté fixe le générateur scalaire \(P'_0=+1\). La condition fixe son signe, et pas seulement sa grandeur. En écrivant
\[
 P_t=\sum_{n\ge0}a_nt^n,
\]
la normalisation détermine \(a_0=a_1=1\). La comparaison du coefficient de \(s^nt\) dans \(P_{s+t}=P_sP_t\) donne \((n+1)a_{n+1}=a_na_1\) ; de manière équivalente, \(P'_t=P_t\). Ainsi,
\begin{equation}
 (n+1)a_{n+1}=a_n,\qquad a_n=\frac1{n!}.
 \label{x_reciprocidad:eq:propagacion}
\end{equation}
Le caractère en une unité de paramètre est \(\Pi_{\rm pr}=P_1\), reconnu ultérieurement comme \(e\).

La normalisation du générateur importe. La loi de semi-groupe sans cette donnée admettrait \(P_t=e^{ct}\) pour différents \(c\) ; c'est pourquoi on n'omet pas la condition qui fixe l'unité du paramètre dans le lecteur. Cette condition est antérieure à la comparaison décimale et appartient à l'exposition de la règle opératoire, et non au choix d'une valeur expérimentale.

\begin{proposition}[Bornes rationnelles de propagation]
Pour \(n\ge1\), soit \(E_n=\sum_{j=0}^n1/j!\). Alors
\[
 E_n<\Pi_{\rm pr}<E_n+\frac{n+2}{(n+1)(n+1)!}.
\]
\end{proposition}
\begin{proof}
Le premier terme omis est \(1/(n+1)!\). Chaque quotient ultérieur est au plus égal à \(1/(n+2)\). La somme géométrique majorant la queue est
\[
 \frac1{(n+1)!}\frac1{1-1/(n+2)}
 =\frac{n+2}{(n+1)(n+1)!}.
\]
L'inégalité stricte découle de la décroissance des quotients successifs.
\end{proof}

\subsection{Incidence modale et génération de l'auto-échelle}

Le sélecteur tritique produit le cycle \((+1,-1,0)\). Dans les modes \(\Sigma,\Pi\), la règle est \(+1:m\mapsto\Sigma\), \(-1:m\mapsto\Pi\), \(0:m\mapsto m\). Avec fermeture cyclique, le mode observé est \((\Sigma,\Pi,\Pi)\). Ses arêtes sont \(\Sigma\to\Pi\), \(\Pi\to\Pi\) et \(\Pi\to\Sigma\). Son incidence, dans cet ordre des modes, est donc
\begin{equation}
 F_{\rm av}=\begin{pmatrix}0&1\\1&1\end{pmatrix}.
 \label{x_reciprocidad:eq:fav}
\end{equation}
Les quatre entrées s'obtiennent à partir du calendrier, et non d'une équation qui aurait préalablement reçu la valeur de \(\varphi\). La répétition du calendrier donne les dénombrements \(3F_{\rm av},9F_{\rm av},36F_{\rm av}\) aux longueurs neuf, vingt-sept et cent huit. Ces dénombrements ne sont pas les puissances de la matrice.

\begin{theorem}[Caractère positif d'auto-échelle]
L'incidence \eqref{x_reciprocidad:eq:fav} possède un unique rayon propre strictement positif. En le normalisant sous la forme \((1,x)^{\mathsf T}\), sa coordonnée \(x\) satisfait \(x^2-x-1=0\), \(x>0\), et vaut \(\varphi=(1+\sqrt5)/2\).
\end{theorem}
\begin{proof}
L'équation propre donne \(x=\lambda\) et \(1+x=\lambda x\). L'élimination de \(\lambda\) produit le polynôme. Une seule de ses deux racines est positive. Comme \(F_{\rm av}^2\) a toutes ses entrées positives, le rayon positif est unique et simple. L'expression radicale reconnaît la coordonnée obtenue ; elle ne participe pas à la production de \(F_{\rm av}\).
\end{proof}

Avec \(F_0=0,F_1=1,F_{n+1}=F_n+F_{n-1}\), l'incidence induit, pour \(n\geq1\),
\[
 F_{\rm av}^n=\begin{pmatrix}F_{n-1}&F_n\\F_n&F_{n+1}\end{pmatrix}.
\]
Les quotients consécutifs \(F_{n+1}/F_n\) alternent autour de \(\varphi\). L'identité de déterminant
\[
 F_{n+1}^2-F_nF_{n+2}=(-1)^n
\]
donne une largeur exacte \(1/(F_nF_{n+1})\) entre deux quotients consécutifs. Sa croissance rend cet encadrement rationnel arbitrairement petit.

La mesure sur toutes les histoires compatibles ne coïncide pas avec la fréquence des modes du calendrier à trois événements. La fréquence chronologique est \((1/3,2/3)\). La mesure d'entropie maximale de l'enveloppe symbolique définie par \(F_{\rm av}\) se construit à partir de ses vecteurs propres et a pour rapport de poids \(\varphi^{-2}\). Cette dernière intervient dans la lecture surfacique--volumétrique ; elle ne s'obtient pas en remplaçant simplement une fréquence de comptage par une autre.


% BEGIN OWNER /Users/ruben/Documents/New project/output/EXTREMA_MEDIA_RAZON_RECIPROCIDAD_ES_EN_20260918/ARTICULO/ES/source/sections/retorno_areal_volumetrico.tex
\subsubsection{Espace des trajectoires et mesure invariante entre feuillets}
\label{x_reciprocidad:sec:medida-retorno-hojas}

La topologie toroïdale d'APP fixe le support des translations cardinales
et leurs classes d'enroulement. La correspondance entre feuillets procède d'
une autre opération sur ce même support : la sélection temporelle de l'
évaluation additive ou multiplicative. Son typage géométrique est
\[
 \eta_{\rm av}(\Sigma)=\mathscr H_{\rm ar},\qquad
 \eta_{\rm av}(\Pi)=\mathscr H_{\rm vol}.
\]
Le premier feuillet correspond à la lecture additive de frontière ; le second,
à la lecture multiplicative d'intérieur. Cette application transporte les étiquettes
de mode et conserve l'incidence \eqref{x_reciprocidad:eq:fav}. Le carré du rapport
d'auto-échelle se détermine au moyen des trajectoires de cette incidence.

\Needspace{8\baselineskip}
\begin{definition}[Enveloppe symbolique à mémoire un]
Soit
\[
 X_{\rm av}=\{(s_k)_{k\in\mathbb Z}\in
 \{\mathscr H_{\rm ar},\mathscr H_{\rm vol}\}^{\mathbb Z}:
 (F_{\rm av})_{s_k,s_{k+1}}=1\text{ pour tout }k\}.
\]
Le décalage des indices agit sur cet espace. C'est le plus petit système
symbolique à mémoire un qui contient les suites modales et conserve
leurs couples consécutifs : il admet les deux transitions croisées et l'
auto-rétention volumétrique. La phase, l'orientation cardinale et les registres
de la trajectoire restent dans l'état TPK à partir duquel on obtient cette
enveloppe.
\end{definition}

\begin{proposition}[Pondération des trajectoires entre feuillets]
Dans l'ordre surfacique--volumétrique, la mesure invariante d'entropie maximale
de \(X_{\rm av}\) a pour matrice de transition et distribution stationnaire
\begin{equation}
 P_{\rm av}=\begin{pmatrix}0&1\\\varphi^{-2}&\varphi^{-1}\end{pmatrix},
 \qquad
 (\mu_{\rm ar},\mu_{\rm vol})
 =\frac{(1,\varphi^2)}{1+\varphi^2}.
 \label{x_reciprocidad:eq:medida-hojas-exacta}
\end{equation}
En particulier, \(\mu_{\rm ar}/\mu_{\rm vol}=\varphi^{-2}\).
\end{proposition}
\begin{proof}
L'incidence déjà construite possède le vecteur propre positif \(r=(1,\varphi)^{\mathsf T}\) et la valeur propre \(\varphi\). La normalisation
\[
 (P_{\rm av})_{ij}
 =\frac{(F_{\rm av})_{ij}r_j}{\varphi r_i}
\]
donne la matrice indiquée. Ses lignes ont pour somme un par \(\varphi^{-2}+\varphi^{-1}=1\). La symétrie de \(F_{\rm av}\) implique que les poids proportionnels à \(r_i^2\) satisfont l'équilibre détaillé ; ainsi, \(\mu P_{\rm av}=\mu\). L'irréductibilité et l'auto-rétention volumique déterminent une unique distribution stationnaire.

Pour vérifier la maximalité, toute transition permise satisfait
\(\log(P_{\rm av})_{ij}=\log r_j-\log r_i-\log\varphi\).
Sous toute mesure invariante, les termes d'extrémité se compensent
en moyenne. Si \(\nu_i\) sont ses poids de sommet et
\(q_i(j)=\nu(s_1=j\mid s_0=i)\), pour \(\nu_i>0\), alors
\[
 h_\nu\leq H_\nu(s_1\mid s_0)
 =\log\varphi-\sum_{i:\nu_i>0}\nu_iD(q_i\Vert P_i)
 \leq\log\varphi,
\]
où \(D(q\Vert p)=\sum_jq_j\log(q_j/p_j)\) est l'entropie
relative sur les couples permis, et l'on adopte \(0\log0=0\).
La mesure markovienne de \eqref{x_reciprocidad:eq:medida-hojas-exacta} atteint la borne.
L'égalité dans la première inégalité exige la propriété de Markov ;
l'égalité dans la seconde exige \(q_i=P_i\). La stationnarité et
l'irréductibilité fixent alors \(\mu\), ce qui établit l'
unicité. C'est la mesure de Parry de l'enveloppe définie.
\end{proof}

La fréquence \((1/3,2/3)\) compte les instants du calendrier primitif ;
\(\mu\) pondère les trajectoires de l'enveloppe. Son rapport quadratique
exprime le produit des amplitudes d'incidence d'entrée et de sortie.
La mémoire et les restrictions de phase du TPK conservent leur domaine
propre ; l'entropie de l'enveloppe correspond au système symbolique
spécifié dans cette définition.

\subsubsection{Opérateur de clôture et limite des retours marqués}
\label{x_reciprocidad:sec:operador-retorno-marcado}

La coordonnée de clôture \(\pi\), générée par le lecteur régional, intervient après la construction de l'incidence modale. Soient \(e_{\rm ar}=(1,0)^{\mathsf T}\), \(e_{\rm vol}=(0,1)^{\mathsf T}\) et \(P_i=e_ie_i^{\mathsf T}\) les projecteurs de feuillet. Les conditions du lecteur sont la conservation des deux feuillets, une normalisation volumique unitaire et une marque surfacique égale à la coordonnée de clôture. Ces conditions déterminent l'opérateur positif
\begin{equation}
 D_\pi=\pi P_{\rm ar}+P_{\rm vol}
       =\begin{pmatrix}\pi&0\\0&1\end{pmatrix}.
 \label{x_reciprocidad:eq:operador-marca-pi}
\end{equation}
La diagonalité exprime la conservation des feuillets et les deux entrées
expriment leurs normalisations ; l'unicité de
l'opérateur pour le lecteur fixé est ainsi explicite.

\Needspace{8\baselineskip}
\begin{proposition}[Rapport des retours de frontière]
Pour \(n\geq1\), l'évaluation marquée
\begin{equation}
 \Theta_n=
 \frac{\langle e_{\rm ar},D_\pi F_{\rm av}^{\,n}e_{\rm ar}\rangle}
 {\langle e_{\rm vol},F_{\rm av}^{\,n}e_{\rm vol}\rangle}
 =\pi\frac{F_{n-1}}{F_{n+1}}
 \label{x_reciprocidad:eq:retorno-marcado-finito}
\end{equation}
satisfait
\begin{equation}
 \lim_{n\to\infty}\Theta_n
 =\pi\frac{\mu_{\rm ar}}{\mu_{\rm vol}}
 =\frac{\pi}{\varphi^2}.
 \label{x_reciprocidad:eq:retorno-marcado-limite}
\end{equation}
\end{proposition}
\begin{proof}
Les entrées diagonales de \(F_{\rm av}^{\,n}\) comptent les chemins
de longueur \(n\) qui reviennent à leur feuillet initial. La formule des puissances
déjà démontrée fournit \(F_{n-1}\) et \(F_{n+1}\). L'opérateur
\(D_\pi\) multiplie par \(\pi\) le premier dénombrement. Enfin,
\(F_{n-1}/F_{n+1}\) est le produit de deux quotients consécutifs
qui convergent vers \(\varphi^{-1}\) ; sa limite est
\(\varphi^{-2}\). L'identité avec le rapport des poids découle de
\eqref{x_reciprocidad:eq:medida-hojas-exacta}.
\end{proof}

L'opérateur de clôture n'apparaît qu'une seule fois, dans l'évaluation terminale du retour. Le transfert \(F_{\rm av}\) reste déterminé par le calendrier. Ainsi se compose l'information de deux lecteurs du même état : l'incidence modale détermine la pondération quadratique et la région de clôture apporte sa coordonnée à la frontière.

La composition conserve aussi le raffinement. Pour des cylindres positifs
\(I=[a,b]\) et \(J=[c,d]\), l'image du lecteur
\(\mathcal R(x,y)=x/y^2\) est l'intervalle
\[
 \mathcal Q(I,J)=\left[\frac{a}{d^2},\frac{b}{c^2}\right].
\]
La monotonie en chaque argument prouve que les cylindres emboîtés de
clôture et d'auto-échelle produisent des intervalles emboîtés du rapport ; la
positivité de leurs limites et la contraction de leurs diamètres déterminent
\(\pi/\varphi^2\). L'expression scalaire retient ainsi une chaîne
d'opérations compatible avec la profondeur, qui s'applique
ci-dessous dans ses deux orientations.

% END OWNER


Une fois les caractères générés, les deux orientations du lecteur
\[
 R(x,y)=\frac{x}{y^2}
\]
fournissent
\begin{equation}
 \rho_A=\frac\pi{\varphi^2},\qquad
 \rho_E=\frac\varphi{\pi^2},\qquad
 \varphi^3\rho_A^2\rho_E=1.
 \label{x_reciprocidad:eq:areal-electron}
\end{equation}
La seconde expression intervient dans l’exposant de la construction électronique de la section~\ref{x_reciprocidad:el-centro:registros}. La première s’obtient également comme limite de la lecture marquée \(\pi F_{n-1}/F_{n+1}\). L’échange des arguments les relie, mais ne transforme pas une reparamétrisation de l’équation électronique en une seconde dérivation indépendante de tous ses coefficients.


% BEGIN OWNER /Users/ruben/Documents/New project/output/EXTREMA_MEDIA_RAZON_RECIPROCIDAD_ES_EN_20260918/ARTICULO/ES/source/sections/euler_trit_articulacion.tex
% Articulación posterior a la generación de pi, phi y e.
% Contenido acumulativo; no reemplaza el desarrollo previo del TRIT.
\subsection{Réalisations exponentielles des coordonnées générées}
\label{x_reciprocidad:euler-trit-articulacion}

Les coordonnées régionales déjà déterminées permettent de reconnaître les
réalisations concrètes de la famille quadratique du TRIT. L’ordre est
substantiel : la construction des régions précède l’évaluation d’une
exponentielle aux paramètres \(\pi\) ou \(2\log\varphi\). Ces expressions
identifient la fonction des coordonnées dans des opérateurs construits à partir
d’APP et du TPK ; elles n’interviennent pas dans la sélection des préfixes qui les produisent.

\subsubsection{Cycle ternaire, transport nilpotent et incidence des feuillets}

La multiplication \(a(r)=4r\) dans \(\mathbb Z/9\mathbb Z\) est d’ordre
\(3\). Sur les unités, elle détermine les orbites \(\{1,4,7\}\) et
\(\{2,8,5\}\). Dans le module d’augmentation de l’une quelconque d’entre elles, formé
des combinaisons entières dont les coefficients ont une somme nulle, la base
\((e_r-e_{4r},e_{4r}-e_{7r})\) donne
\[
 C=\begin{pmatrix}0&-1\\1&-1\end{pmatrix},
 \qquad C^2+C+I=0.
\]
La forme restreinte a pour matrice de Gram
\(\bigl(\begin{smallmatrix}2&-1\\-1&2\end{smallmatrix}\bigr)\).
Dans les évaluations résiduelles modulo \(9\), l’action distingue les feuillets :
\(S(ax,ay)=aS(x,y)\), tandis que \(P(ax,ay)=a^{-1}P(x,y)\).
La réalisation cyclique conserve donc une orientation contragrédiente
entre somme et produit. Les quotients du relèvement entier sont conservés
dans le registre arithmétique correspondant.

Le transport parabolique élémentaire est représenté par
\[
 N=\begin{pmatrix}0&1\\0&0\end{pmatrix}.
\]
L’incidence surfacique--volumique déjà construite dans \eqref{x_reciprocidad:eq:fav} fournit
\[
 F_{\rm av}=\begin{pmatrix}0&1\\1&1\end{pmatrix},
 \qquad A=F_{\rm av}^2=\begin{pmatrix}1&1\\1&2\end{pmatrix}.
\]
Les trois opérateurs agissent dans des espaces réels déclarés : le module
d’augmentation tensorisé avec \(\mathbb R\), le plan du transport nilpotent et
le plan d’incidence modale. Le partage d’une dimension matricielle n’identifie
ni ces espaces ni leurs coordonnées.

\begin{proposition}[Générateurs concrets des trois types]
\label{x_reciprocidad:euler-trit-generadores}
Les opérateurs
\[
 J_{\mathrm e}=\frac{2C+I}{\sqrt3},\qquad
 J_{\mathrm p}=N,\qquad
 J_{\mathrm h}=\frac{2A-3I}{\sqrt5}
\]
satisfont, respectivement,
\[
 J_{\mathrm e}^{\,2}=-I,\qquad
 J_{\mathrm p}^{\,2}=0,\qquad
 J_{\mathrm h}^{\,2}=I.
\]
\end{proposition}
\begin{proof}
La relation de \(C\) implique \((2C+I)^2=-3I\).
La multiplication directe donne \(N^2=0\).
Enfin, \(A\) est de trace \(3\) et de déterminant \(1\), d’où
\(A^2-3A+I=0\) et \((2A-3I)^2=5I\).
\end{proof}

\begin{proposition}[Trois évaluations de l’identité exponentielle]
\label{x_reciprocidad:euler-trit-evaluaciones}
Dans les réalisations précédentes,
\begin{equation}
 \boxed{
 \exp(\pi J_{\mathrm e})+I=0,\qquad
 \exp(J_{\mathrm p})=I+J_{\mathrm p},\qquad
 \exp(2\log\varphi\,J_{\mathrm h})=F_{\rm av}^{\,2}.}
 \label{x_reciprocidad:euler-trit-tres-evaluaciones}
\end{equation}
\end{proposition}
\begin{proof}
La spécialisation elliptique de l’exponentielle tritique donne
\(\cos\pi\,I+\sin\pi\,J_{\mathrm e}=-I\).
La nilpotence élimine exactement les termes de degré au moins \(2\)
dans la deuxième expression. Pour la troisième, \(\varphi^2=\varphi+1\) implique
\[
 \cosh(2\log\varphi)=\frac32,\qquad
 \sinh(2\log\varphi)=\frac{\sqrt5}{2}.
\]
L’exponentielle vaut donc
\(\frac32I+\frac{\sqrt5}{2}J_{\mathrm h}=A\).
\end{proof}

La première égalité représente le demi-tour elliptique ; le tour complet
satisfait \(\exp(2\pi J_{\mathrm e})=I\). La deuxième maintient un terme
linéaire non nul : le régime parabolique exprime un transport, non une absence
d’évolution. La troisième réalise une progression hyperbolique unimodulaire. Ses
valeurs propres sont \(\varphi^2\) et \(\varphi^{-2}\) ; une direction se dilate et
l’autre se contracte en conservant le déterminant. Les trois évaluations
emploient des paramètres distincts d’une famille exponentielle commune.

Dans cette famille, \(e\) intervient par l’évaluation scalaire
\(\exp(I)=eI\) et la loi \(\exp(sI)\exp(tI)=\exp((s+t)I)\).
Son rôle ne se restreint pas au type parabolique. La coordonnée \(e\), obtenue
auparavant à partir de sa région, est reconnue ici comme l’évaluation unitaire
de l’exponentielle ; cette identité ne remplace pas sa généalogie coinductive.

\subsubsection{Résolution polynomiale des régimes}

Les trois réalisations admettent une présentation compacte sur la somme
directe de leurs espaces. On définit
\[
 \mathbb J=J_{\mathrm e}\oplus J_{\mathrm p}\oplus J_{\mathrm h},
 \qquad Q=\mathbb J^2.
\]
Les relations précédentes impliquent
\[
 \mathbb J^6=\mathbb J^2,\qquad Q^3=Q.
\]
L’emploi de \(Q\) maintient ce carré opératoriel distinct du registre
dodécaphasique \(K\).

\begin{proposition}[Projecteurs polynomiaux de régime]
\label{x_reciprocidad:euler-trit-proyectores}
Les applications
\[
 p_{\mathrm e}=\frac{Q^2-Q}{2},\qquad
 p_{\mathrm p}=I-Q^2,\qquad
 p_{\mathrm h}=\frac{Q^2+Q}{2}
\]
sont des idempotents mutuellement orthogonaux, de somme \(I\), qui restituent les trois
termes de la somme directe représentant les régimes.
\end{proposition}
\begin{proof}
Les trois polynômes prennent les valeurs \((1,0,0)\), \((0,1,0)\) et
\((0,0,1)\) aux points \(-1,0,+1\), respectivement. Les différences
entre leurs carrés et eux-mêmes, ainsi que leurs produits croisés, sont
des multiples de \(X^3-X\). Substituer \(Q\), qui annule ce polynôme, démontre
les identités. Dans les trois blocs, \(Q\) agit comme \(-I,0,I\).
\end{proof}

On obtient ainsi
\begin{align}
 \exp(t\mathbb J)
 ={}&p_{\mathrm e}(\cos t\,I+\sin t\,\mathbb J)
   +p_{\mathrm p}(I+t\mathbb J)\notag\\
   &+p_{\mathrm h}(\cosh t\,I+\sinh t\,\mathbb J).
 \label{x_reciprocidad:euler-trit-exponencial-total}
\end{align}
Chaque terme conserve sa relation quadratique et la conjugaison inverse
son évolution. La représentation totale réunit les trois régimes, mais ne
définit pas à elle seule un opérateur de transition entre eux. Une telle transition
requiert les applications de transport du TPK avec leurs domaines, leur orientation et
leur mémoire. L’articulation exponentielle permet de reconnaître leurs types sans
remplacer cette dynamique par une identification des fibres.

% PROCEDENCIA FOCAL (no impresa)
% Resultado recuperado: tres generadores y tres evaluaciones exponenciales.
% Propietario completo:
% BASE/manuscrito/sections/hmt/09b_algebras_cuadraticas.tex:289-516.
% Antecedente causal de C y acción contragrediente:
% BASE/manuscrito/sections/hmt/03d_esqueleto_ciclotomico_app.tex:13-47,145-246.
% Los tres tipos y la orientación están en:
% BASE/manuscrito/sections/hmt/02_trit_geometrias.tex:83-196.
% Resolución polinómica recuperada:
% output/INVESTIGACION_HOLONOMIA_NONADICA_CRISTAL_APERIODICO_20260903/
% ALGEBRA_HOLONOMICA_UNIVERSAL_Y_EULER_TRITICO.md, secciones 8-11 y 22.
% Los cuatro propietarios se leyeron completos. BASE designa:
% /Users/ruben/Documents/HMT2/
% EL_CIERRE_HOLOGRAFICO_DEL_INFINITO_SUCESOR_102_CAPITULOS_2026-08-28_EN_TRABAJO/
% No se incorpora la identidad P9(alpha)=delta del documento especializado:
% esa identificación no pertenece a este bloque.
% Tampoco se identifica una dirección nilpotente con la vacancia electrónica,
% ni se presume una representación global del transporte por la suma directa.
% COMPROBACIONES: Python estándar, enteros y Fraction.
% C^2+C+I=0; (2C+I)^2=-3I; N^2=0; (2Fav^2-3I)^2=5I.
% 81 pares de residuos verifican las dos acciones contragredientes.
% Los 3 proyectores se verificaron en Q=-1,0,+1; prueba general en texto.
% 6 labels únicos y entornos equilibrados. Sin compilación.

% END OWNER


\subsection{Cylindres compatibles et profondeur arbitraire}

La réalisation d'un bloc ternaire \(b=(b_1,\ldots,b_6)\) est l'entier
\[
 \nu(b)=\sum_{j=1}^6b_j3^{6-j}\in\{0,\ldots,728\}.
\]
Une suite de blocs définit des préfixes entiers par \(N_{n+1}=729N_n+\nu(b_{n+1})\), en commençant en \(N_0=m\), où \(m\) est la partie entière déterminée par le lecteur. Les fourchettes précédentes situent les caractères de clôture, de propagation et d'auto-échelle dans \((3,4)\), \((2,3)\) et \((1,2)\), respectivement ; leurs valeurs initiales sont donc \(m=3,2,1\). Les blocs suivants raffinent la partie fractionnaire. Le cylindre correspondant est
\begin{equation}
 I_n=\left[\frac{N_n}{729^n},\frac{N_n+1}{729^n}\right],
 \qquad |I_n|=729^{-n}.
 \label{x_reciprocidad:eq:cilindro}
\end{equation}
La convention semi-ouverte est utilisée pour sélectionner un représentant lorsqu'une valeur tombe sur une frontière. Pour le passage à la limite, on utilise les fermetures de ces intervalles et on identifie les deux développements de frontière équivalents. Cette précision empêche de supposer que toute intersection d'intervalles semi-ouverts emboîtés est automatiquement non vide.

\begin{theorem}[Détermination de la réalisation scalaire]
Une histoire de blocs compatible détermine une unique limite des extrémités gauches de \eqref{x_reciprocidad:eq:cilindro}. Si un caractère exact dispose d'encadrements rationnels de largeur arbitrairement petite et se sépare de la frontière de chaque cylindre décimal examiné, chaque préfixe décimal de longueur finie est déterminé après un calcul fini. Les préfixes ainsi obtenus sont compatibles entre eux.
\end{theorem}
\begin{proof}
L'extension d'un bloc conserve le préfixe précédent par division euclidienne. Les cylindres fermés sont emboîtés et leur diamètre tend vers zéro ; leurs extrémités gauches forment une suite de Cauchy et possèdent une unique limite. Pour une longueur décimale fixée, une valeur séparée des frontières de la partition possède une distance positive à celles-ci. Une fourchette suffisamment étroite est contenue dans une unique cellule et détermine son préfixe. L'unicité de la cellule et l'emboîtement des partitions prouvent la compatibilité par troncature. Aux points à développement terminal, il faut ajouter la décision exacte de frontière et sa convention de représentation ; cette décision ne se déduit pas d'une largeur numérique seule.
\end{proof}

\begin{corollary}[Obtention finie et compatibilité des préfixes régionaux]
\label{x_reciprocidad:gen:prefijos-regionales}
Pour chacune des réalisations \(\pi,\varphi,e\) et pour toute longueur
décimale finie, les encadrements construits déterminent le préfixe correspondant
par un calcul fini. Les préfixes obtenus à différentes profondeurs
sont compatibles par troncature.
\end{corollary}
\begin{proof}
Les caractères de clôture, de propagation et d'auto-échelle possèdent les encadrements précédents. Après leur reconnaissance, l'irrationalité de \(\pi,e,\varphi\) garantit leur séparation des frontières rationnelles de chaque partition décimale finie. On applique le théorème précédent à chaque caractère.
\end{proof}
La profondeur arbitraire exprime le quantificateur du corollaire : pour tout
horizon fini, il existe un calcul qui détermine son préfixe. L'objet
mathématique est la famille compatible complète ; chaque exécution en obtient une
projection finie.

\subsection{Involution spéculaire et synchronisation des régions}

La connexion nonadique transporte conjointement les trois caractères. Dans la carte de la neuvième phase, l'axe d'auto-échelle et la somme ternaire de clôture et de propagation restent fixes sous l'échange spéculaire. L'inversion hélicoïdale change le sens du transport ; le miroir des deux canaux change leur orientation relative. Ce sont des involutions différentes, bien qu'elles agissent toutes deux sur le même système de trajectoires.

Les coïncidences de recensement illustrent cette coordination :
\begin{center}
\begin{tabular}{@{}c rrr l@{}}\toprule
Phase&\(N_\pi\)&\(N_e\)&\(N_\varphi\)&Égalité des dénombrements\\\midrule
4&207&207&122&\(N_\pi=N_e\)\\
5&39&151&151&\(N_e=N_\varphi\)\\
6&110&110&69&\(N_\pi=N_e\)\\
7&81&29&81&\(N_\varphi=N_\pi\)\\
8&59&59&59&synchronisation triple\\
9&43&43&43&synchronisation triple\\\bottomrule
\end{tabular}
\end{center}
Ces entiers comptent des coordonnées compatibles des régions. Ils n'affirment pas une égalité entre \(\pi\), \(e\) et \(\varphi\) en tant que nombres réels. Si
\[
 \partial(a,b,c)=(a-b,b-c,c-a),
\]
son image appartient au réseau \(A_2=\{(u,v,w)\in\ZZ^3:u+v+w=0\}\). Aux phases huit et neuf, la composante relative s'annule, mais la composante diagonale passe de \((59,59,59)\) à \((43,43,43)\). La synchronisation relative coexiste avec un déplacement de \(-16\) par canal. La fermeture qui conduira à \(\alpha\) doit conserver les deux informations : coïncidence relative et évolution de la mémoire.

\subsection{Suspension hélicoïdale sturmienne}

La relation entre les bases neuf et dix de la complétion définit le paramètre adimensionnel
\begin{equation}
 \delta=\frac{\log(10/9)}{\log10}.
 \label{x_reciprocidad:eq:delta-sturm}
\end{equation}
Elle ne s'obtient pas à partir d'une constante physique cible. C'est l'échelle logarithmique relative de deux bases déjà présentes dans la construction. Son emploi comme lecture symbolique du raffinement a une conséquence exacte.

\begin{lemma}[Irrationalité de la pente]
Le nombre \(\delta\) est irrationnel.
\end{lemma}
\begin{proof}
Si \(\delta=p/q\), \(q>0\), alors \((10/9)^q=10^p\) et \(9^q=10^{q-p}\). La factorisation du premier membre contient le nombre premier trois et celle du second uniquement les nombres premiers deux et cinq. L'égalité est impossible.
\end{proof}

Le mot mécanique inférieur, d'ordonnée à l'origine \(\theta\), est
\[
 z_n(\theta)=\lfloor(n+1)\delta+\theta\rfloor-\lfloor n\delta+\theta\rfloor\in\{0,1\}.
\]
On l'obtient en observant la rotation \(\theta\mapsto\theta+\delta\pmod1\) au moyen de l'intervalle \([1-\delta,1)\). La suite d'événements a pour densité \(\delta\), et la somme sur toute fenêtre de longueur \(L\) vaut \(\lfloor L\delta\rfloor\) ou \(\lceil L\delta\rceil\). Puisque \(21<1/\delta<22\), les séparations entre événements sont vingt et un ou vingt-deux. Le mot qui indique quels sont les intervalles courts est un autre mot mécanique de pente \(\sigma=22-1/\delta\) ; ses événements sont séparés par six ou sept intervalles primaires. La seconde échelle est induite par la première et n'est pas un second paramètre libre.


% BEGIN OWNER /Users/ruben/Documents/New project/output/EXTREMA_MEDIA_RAZON_RECIPROCIDAD_ES_EN_20260918/ARTICULO/ES/source/sections/vacancias_capacidad.tex
% Desarrollo reunido de DESARROLLO_MATEMATICO, apartados 38--40 y 50.
\subsubsection{Indice de capacité, lacunes et induction du régime tritique}
\label{x_reciprocidad:vac-capacidad-trit}

La suite des lacunes est liée au raffinement par une
identité d'incréments. Pour la montrer, on fixe l'origine de phase
\(\theta=0\). Soient
\[
 \rho=\log_{10}9=1-\delta,\qquad
 \mathcal C_t=\lfloor(t+1)\rho\rfloor\quad(t\geq0).
\]
L'indice \(\mathcal C_t\) exprime la capacité décimale du raffinement
dans la normalisation qui compare un bloc ternaire de \(6\) positions,
avec \(729\) possibilités, et un bloc décimal de \(3\) positions, avec
\(1000\) possibilités. En effet,
\(\log_{1000}729=\log_{10}9=\rho\). La lettre \(K\) est réservée au
registre dodécaphasique qui sera construit ultérieurement.

\begin{proposition}[Complémentarité de la capacité et de la lacune]
\label{x_reciprocidad:prop:capacidad-vacancia}
Pour \(t\geq1\),
\begin{equation}
 \mathcal C_t-\mathcal C_{t-1}=1-z_t(0).
 \label{x_reciprocidad:eq:capacidad-vacancia}
\end{equation}
Par conséquent, si \(0\leq a<b\),
\[
 \mathcal C_b-\mathcal C_a+
 \sum_{t=a+1}^{b}z_t(0)=b-a.
\]
\end{proposition}
\begin{proof}
L'irrationalité de \(\delta\) donne
\[
 \mathcal C_t
 =t+1-\lceil(t+1)\delta\rceil
 =t-\lfloor(t+1)\delta\rfloor.
\]
La soustraction de deux termes consécutifs produit
\eqref{x_reciprocidad:eq:capacidad-vacancia}. La somme sur l'intervalle est télescopique.
\end{proof}

Chaque transition augmente d'une unité le compteur des blocs. Si
\(z_t=0\), la capacité augmente également ; si \(z_t=1\), la capacité se
maintient tandis qu'une transition est incorporée à l'histoire. La lacune
est ici un événement de rétention de capacité, de position et
d'antécédent déterminés. Le bilan précédent conserve exactement la
longueur entre capacité acquise et lacunes enregistrées. La phase
observée \(g_t^{\rm cap}=[\mathcal C_t]_9\) vérifie
\[
 g_t^{\rm cap}-g_{t-1}^{\rm cap}\equiv1-z_t\pmod9.
\]
Cette phase de capacité et la phase chronologique \([t]_9\) sont deux
observations différentes. Une rétention locale de capacité et un tour
complet de l'holonomie ont des mécanismes distincts, bien que les deux
permettent de conserver l'observation de phase avec modification de l'état.

La relation avec TRIT peut se suivre dans les positions mêmes des lacunes.
Pour \(j\geq1\), définissons
\[
 p_j=\left\lfloor\frac j\delta\right\rfloor,\qquad
 v_j=\mathcal C_{p_j},\qquad h_j=p_{j+1}-p_j.
\]
L'inégalité \(p_j\delta<j<(p_j+1)\delta\) identifie \(p_j\)
comme la position de l'événement \(j\). Comme \(0<\delta<1\), elle donne aussi
\(\lfloor(p_j+1)\delta\rfloor=j\), et par conséquent
\[
 v_j=p_j-j,\qquad
 v_{j+1}-v_j=h_j-1.
\]
Si \(s_j=22-h_j\), alors \(s_j\in\{0,1\}\) et
\begin{equation}
 [v_{j+1}]_3=[v_j-s_j]_3.
 \label{x_reciprocidad:eq:vacancia-regimen}
\end{equation}
Un intervalle de longueur \(22\) conserve la classe ; un intervalle de longueur
\(21\) la décale d'une unité dans le sens négatif. Le sélecteur canonique
s'applique à la phase positive \(r_j=1+[v_j]_9\) :
\[
 \tau_j=M_{\rm ph}(r_j).
\]
Un régime tritique est ainsi déterminé par l'indice de capacité
retenu à chaque événement, en plus du codage binaire de la lacune.

Les positions des intervalles courts forment le codage mécanique
induit de pente \(22-1/\delta\). Leurs espacements \(6\) et \(7\)
déterminent les longueurs des plateaux de la classe \([v_j]_3\), une
fois l'origine et la convention des extrémités fixées. Les premières classes
sont \(2\) pendant \(6\) événements, \(1\) pendant \(7\) et \(0\) pendant
\(7\) ; le premier segment de \(6\) est la restriction initiale d'un plateau.
Leurs signatures parcourent \(0,-1,+1\). Les équations
\eqref{x_reciprocidad:eq:capacidad-vacancia} et \eqref{x_reciprocidad:eq:vacancia-regimen} explicitent le
lien entre raffinement, lacunes et sélection du régime. La lecture
quadratique de TRIT réalise ensuite ces types au moyen de leurs générateurs
respectifs ; l'événement temporel et sa représentation géométrique
conservent ainsi une dépendance identifiable.

L'inversion de l'indice monotone fournit une autre observation utile :
\[
 \mathcal T(k)=\min\{t:\mathcal C_t=k\}
   =\left\lceil\frac k\rho\right\rceil-1\qquad(k\geq1).
\]
Si \(\sigma=\rho^{-1}-1\), le nombre de blocs nécessaires pour acquérir
\(m\) unités de capacité est
\[
 \mathcal T(k+m)-\mathcal T(k)
 =m+\lceil(k+m)\sigma\rceil-\lceil k\sigma\rceil
 \in\{m+\lfloor m\sigma\rfloor,m+\lceil m\sigma\rceil\}.
\]
En particulier, \(9\) unités de capacité nécessitent \(9\) ou \(10\)
blocs. La période d'une phase finie reste exacte, tandis que le
temps de retour mesuré par l'autre indice admet deux valeurs. Cette
distinction fait partie de la dynamique du raffinement et ne constitue pas une correction
de ses valeurs numériques a posteriori.

% END OWNER


\begin{theorem}[Complexité des lectures de phase]
Pour \(m\ge1\), le mot mécanique a une complexité factorielle \(p(m)=m+1\). En conservant la phase \(n\bmod9\), la complexité est \(p_9(m)=9(m+1)\) ; en conservant seulement sa classe ternaire, elle est \(p_3(m)=3(m+1)\). Les trois entropies topologiques sont nulles.
\end{theorem}
\begin{proof}
Les prédécesseurs des deux extrémités de l'intervalle de codage forment, pour les facteurs de longueur \(m\), \(m+1\) points distincts du cercle : la coïncidence entre extrémités translatées élimine les doublons et l'irrationalité empêche toute coïncidence supplémentaire. Ces points déterminent \(m+1\) intervalles à itinéraires constants et distincts. Toute orbite irrationnelle visite chaque intervalle, d'où \(p(m)=m+1\).

Une phase modulo neuf étant fixée, les positions de cette phase parcourent une rotation de pas \(9\delta\), également irrationnel. Par densité, tous les facteurs apparaissent dans chaque phase, et la première coordonnée distingue leurs neuf copies. Le même argument avec le pas \(3\delta\) donne trois copies après réduction ternaire. Enfin, \(m^{-1}\log(c(m+1))\to0\) pour \(c=1,3,9\).
\end{proof}

Le produit de phases \(C_9\times\mathbb T\), translaté par \((1,\delta)\), possède des caractères de fréquence
\[
 \frac a9+k\delta\pmod1,\qquad a\in\ZZ/9\ZZ,\ k\in\ZZ.
\]
Le module spectral n'est pas \(\log10\) : c'est le module additif engendré par \(1/9\) et \(\delta\), modulo les entiers. Le facteur \(\log10\) appartient à la normalisation de \eqref{x_reciprocidad:eq:delta-sturm}.

Une présentation opératorielle rend la mémoire explicite. Dans \(\ell^2(\ZZ)\), soient \(T|n\rangle=|n+1\rangle\), \(V_9|n\rangle=\zeta_9^n|n\rangle\), \(W_\delta|n\rangle=e^{2\pi i n\delta}|n\rangle\) et \(D_M|n\rangle=\lfloor n/9\rfloor|n\rangle\). Sur les vecteurs à support fini,
\begin{equation}
 V_9T=\zeta_9TV_9,\qquad
 W_\delta T=e^{2\pi i\delta}TW_\delta,\qquad
 [D_M,T^9]=T^9.
 \label{x_reciprocidad:eq:weyl-relojes}
\end{equation}
Les deux premières relations sont des covariances de Weyl ; la troisième exprime qu'un tour nonadique augmente le degré de mémoire. La rotation de base est abélienne, tandis que l'algèbre des observables et des translations ne l'est pas. L'inversion \(|n\rangle\mapsto|-n\rangle\) conjugue l'avancée avec le recul. Cette réalisation fournit une suspension hélicoïdale apériodique à retour fini de phase, et non un mot périodique de longueur neuf.


% BEGIN OWNER /Users/ruben/Documents/New project/output/EXTREMA_MEDIA_RAZON_RECIPROCIDAD_ES_EN_20260918/ARTICULO/ES/source/figures/holonomia.tex
\begin{figure}[htbp]
\centering
\begin{tikzpicture}[x=1cm,y=1cm,>=Latex,font=\small]
\begin{scope}[xshift=1.8cm,yshift=.3cm]
\draw[->,gray] (0,-.5)--(0,3.65) node[above] {\(n/9\)};
\foreach \k in {0,1,2,3}{
 \draw[gray!60,dashed] (-1.1,\k)--(1.1,\k);
 \node[left,font=\scriptsize] at (-1.1,\k){\(\k\)};
}
\draw[black,line width=.65pt,domain=0:27,samples=260,smooth,variable=\n]
 plot ({cos(40*\n)},{\n/9+.14*sin(40*\n)});
\foreach \n in {0,...,27}{
 \fill ({cos(40*\n)},{\n/9+.14*sin(40*\n)}) circle (.026);
}
\foreach \k in {0,1,2,3}{\fill (1,\k) circle (.045);}
\node[align=center,text width=3.3cm,font=\scriptsize] at (0,-1.02)
 {Même phase en \(n=0,9,18,27\) ;\\degré de mémoire distinct.};
\end{scope}
\begin{scope}[xshift=4.2cm,yshift=.5cm]
\node[anchor=west] at (0,3.1){\(z_n(0)=\lfloor(n+1)\delta\rfloor-\lfloor n\delta\rfloor\)};
\draw[->] (0,1.7)--(9,1.7) node[right]{\(n\)};
\foreach \n in {21,43,65,87,109,131,152,174,196,218}{
 \draw[line width=.6pt] ({\n*.039},1.7)--({\n*.039},2.25);
 \fill ({\n*.039},2.25) circle (.033);
 \node[below,font=\tiny] at ({\n*.039},1.7){\n};
}
\foreach \a/\b/\g in {21/43/22,43/65/22,65/87/22,87/109/22,109/131/22,131/152/21,152/174/22,174/196/22,196/218/22}{
 \node[font=\scriptsize] at ({(\a+\b)*.0195},2.55){\g};
}
\node[anchor=west,align=left,text width=8.3cm] at (0,.55)
 {Le codage irrationnel détermine les intervalles entre événements. Le retour de la phase nonadique ne transforme pas cette suite en un mot périodique.};
\end{scope}
\end{tikzpicture}
\caption{Représentations complémentaires du transport. À gauche, projection de l'hélice \((\cos(2\pi n/9),\sin(2\pi n/9),n/9)\) : l'avancée de neuf pas retrouve la phase et augmente la hauteur. À droite, les dix premiers événements du mot inférieur pour \(\delta=\log_{10}(10/9)\), avec des espacements exacts de \(21\) ou \(22\). La représentation hélicoïdale n'introduit pas de métrique physique supplémentaire.}
\label{x_reciprocidad:fig:holonomia-esturmiana}
\end{figure}

% END OWNER


Pour ce dénombrement, nous fixons l'interception \(\theta=0\) et les extrémités cumulées \((N_0,\ldots,N_4)=(0,729,972,999,1000)\), correspondant à l'ordre cubique \((729,243,27,1)\). Les dénombrements inférieurs sont \(\lfloor N_j\delta\rfloor-\lfloor N_{j-1}\delta\rfloor\) et donnent \((33,11,1,0)\). Les supérieurs sont \(\lceil N_j\delta\rceil-\lceil N_{j-1}\delta\rceil\), avec la même frontière initiale nulle, et donnent \((34,11,1,0)\). Ces valeurs dépendent de l'interception et de l'ordre des strates ; elles ne sont pas affirmées pour une phase initiale arbitraire. Leurs totaux sont quarante-cinq et quarante-six. L'incrément correspond à une unique unité de frontière, située dans la première strate de cette partition ordonnée. Les nombres trente-quatre et onze apparaissent ainsi dans un même dénombrement antérieur à toute carte d'unités ; les identifier à des exposants physiques exige en outre la dérivation dimensionnelle correspondante.

La dénomination \emph{cristal apériodique} est employée ici pour cet ordre symbolique, dont la complexité, le spectre et le retour sont spécifiés. L'entropie topologique nulle n'équivaut pas à elle seule à une dissipation physique nulle : une réalisation thermodynamique exige une dynamique énergétique et une opération d'effacement définies. De même, cette lecture sturmienne ne remplace pas tous les opérateurs du TPK et ne démontre pas à elle seule la périodicité ou l'apériodicité de la sortie complète de \(\alpha\). Sa fonction est d'expliciter une structure d'ordre que le simple retour du calendrier ne décrit pas.


% BEGIN OWNER /Users/ruben/Documents/New project/output/EXTREMA_MEDIA_RAZON_RECIPROCIDAD_ES_EN_20260918/ARTICULO/ES/source/sections/revision_monodromia.tex
% Adición acumulativa. Insertar después de la suspensión esturmiana existente.
% Conservar íntegro el texto y la figura fig:holonomia-esturmiana.
\subsection{Connexion nonadique, holonomie et représentation monodromique}
\label{x_reciprocidad:mono-ampl:conexion}

La connexion régionale déjà construite est maintenant étudiée au moyen de trois objets : le transport de l'état complet, la représentation de ses tours et la suspension de ses itinéraires. Les opérations d'APP--TRIT--TPK et le catalogue initial conservent les définitions précédentes. La distinction nécessaire concerne ici ce que chaque représentation retient d'une même histoire (sections~\ref{x_reciprocidad:sec:extension} et~\ref{x_reciprocidad:sec:generacion}).

À l'indice de prolongement \(k\), une carte de l'état conserve
\[
 \widehat x_k=(B_k,C_k,p_k,c_k,\Sigma_k,W_k,g_k,m_k).
\]
Ici, \(B_k\) réunit les trois blocs régionaux, \(C_k\) est le cylindre,
\(p_k\) le préfixe, \(c_k\) le registre de retenue arithmétique, \(\Sigma_k\) la signature
de frontière et \(W_k\) l'information incidencielle. La phase est
\(g_k=1+[(k-1)]_9\) ; \(m_k\) compte les tours complets. La réduction
\(\beta_k=[m_k]_{12}\) conserve la position dodécaphasique, tandis que
l'entier \(m_k\) distingue ses tours successifs.

\begin{definition}[Transport d'un tour nonadique]
\label{x_reciprocidad:mono-ampl:def-vuelta}
Soient \(\mathcal R_k\subseteq\widehat X_k\times\widehat X_{k+1}\)
les relations d'extension admissible du TPK. La réalisation de
l'holonomie au niveau \(k\) est la composition relationnelle
\[
 \Gamma_{9,k}=\mathcal R_{k+8}\circ\cdots\circ\mathcal R_k.
\]
Pour \(s\geq1\), son prolongement à \(s\) tours est
\[
 \Gamma_{9s,k}=
 \Gamma_{9,k+9(s-1)}\circ\cdots\circ\Gamma_{9,k}.
\]
\end{definition}

La formulation relationnelle conserve les continuations admissibles. Sur
une histoire individualisée, chaque flèche mène à l'état effectivement
prolongé ; le bloc visible isolé peut encore admettre plusieurs
continuations. Cette différence est précisément la fonction de la mémoire et
des conditions de frontière dans la détermination de la trajectoire.

\begin{proposition}[Itération du retour et mémoire accumulée]
\label{x_reciprocidad:mono-ampl:vueltas}
Dans toute composition admissible de \(s\) tours,
\[
 g_{k+9s}=g_k,\qquad m_{k+9s}=m_k+s,\qquad
 \beta_{k+9s}=\beta_k+s\pmod{12}.
\]
\end{proposition}
\begin{proof}
La règle d'un tour conserve la phase et incrémente le compteur d'une
unité. Son application à l'état déjà transporté répète exactement ces
deux opérations. La récurrence donne les deux premières égalités ; réduire la
seconde modulo \(12\) donne la troisième.
\end{proof}

Le compteur fournit ainsi un témoin entier de la différence entre les
tours. Après \(12\) tours, le couple des phases finies est rétabli, mais
le compteur augmente de \(12\). L'extension non scindée
\eqref{x_reciprocidad:eq:extension108} exprime la composition du calendrier par son
cocycle ; la trajectoire conserve en outre le relèvement entier de cette donnée.
La conservation signifie ici le transport des contributions antérieures,
et non l'immobilité de toutes les coordonnées.

\subsubsection{Représentation bilatérale de l'indice de retour}

Le décalage \(T\) de \eqref{x_reciprocidad:eq:weyl-relojes} admet une
expression en coordonnées de phase et de tour. Numérotons les phases de \(0\) à
\(8\), par un décalage de l'origine par rapport à \(g_k\), et
écrivons \(k=9m+r\), avec \(0\leq r<9\), sur la droite entière
qui revêt le cycle des phases. Sur
\(\mathcal H_{\rm ind}=\ell^2(\mathbb Z)\), la base
\(|k\rangle=|r,m\rangle\) exprime ce même décalage unitaire
\[
 T|r,m\rangle=
 \begin{cases}
 |r+1,m\rangle,&r<8,\\
 |0,m+1\rangle,&r=8.
 \end{cases}
\]
L'inverse réduit \(r\) d'une unité lorsque \(r>0\), et envoie
\(|0,m\rangle\) sur \(|8,m-1\rangle\). Le passage par l'origine de phase
incorpore donc l'unité exacte de quotient.

\Needspace{8\baselineskip}
\begin{proposition}[Représentation de la monodromie du cycle]
\label{x_reciprocidad:mono-ampl:representacion}
Soit \(S=T^9\). Alors
\[
 S|r,m\rangle=|r,m+1\rangle,\qquad
 \mathbb Z\longrightarrow\mathcal U(\mathcal H_{\rm ind}),
 \quad s\longmapsto S^s
\]
est une représentation unitaire fidèle du groupe des tours du cycle.
\end{proposition}
\begin{proof}
Neuf décalages conservent le reste de la division par \(9\) et
incrémentent le quotient. Les puissances satisfont
\(S^{s+t}=S^sS^t\), et \(S^s|r,m\rangle=|r,m+s\rangle\) distingue
tout entier \(s\). La permutation d'une base orthonormée est unitaire.
\end{proof}

Cette représentation correspond au revêtement du cycle, dont le groupe
des tours est \(\mathbb Z\), et représente la phase avec son compteur. L'état
du TPK contient en outre les coordonnées de trajectoire déjà décrites.
L'extension bilatérale admet des indices négatifs comme coordonnées du
revêtement ; une trajectoire issue d'une condition initiale utilise
son demi-axe admissible. On dispose ainsi d'un inverse opératoriel
de l'indice sans transformer chaque relation d'extension en bijection de
l'état complet.

Avec \(D_M|k\rangle=\lfloor k/9\rfloor|k\rangle\), l'identité
\([D_M,S]=S\) exprime l'incrément de mémoire sur le domaine dense des
vecteurs à support fini. Si \(\mathcal I|k\rangle=|-k\rangle\) et
\(P_0\) projette sur les indices divisibles par \(9\), on obtient
également
\begin{equation}
 \mathcal I T\mathcal I=T^{-1},\qquad
 \mathcal I D_M\mathcal I=-D_M-(I-P_0).
 \label{x_reciprocidad:mono-ampl:inversion-contador}
\end{equation}
En effet, pour \(k=9m+r\),
\(\lfloor-k/9\rfloor=-m\) si \(r=0\) et vaut \(-m-1\) si \(r>0\).
L'inversion conserve ainsi le terme de frontière de la division euclidienne.

\subsection{Horloge des blocs et horloge de capacité}
\label{x_reciprocidad:mono-ampl:relojes}

L'indice d'une application de raffinement et le nombre de positions
décimales disponibles ont des lois liées, mais différentes. Pour
les exprimer, nous réservons \(n\) à l'indice des blocs ternaires et définissons
\[
 \rho=\log_{1000}729=\log_{10}9,\qquad
 \delta=1-\rho,\qquad
 \mathcal C_n=\lfloor(n+1)\rho\rfloor,\quad n\geq0.
\]
La lettre \(\mathcal C_n\) désigne ici la capacité, et non le registre
dodécaphasique \(K\). À l'origine de phase canonique, le mot mécanique
précédent fournit
\[
 Z_n=\sum_{j=0}^{n}z_j(0)=\lfloor(n+1)\delta\rfloor.
\]

\begin{proposition}[Bilan exact entre capacité et lacunes]
\label{x_reciprocidad:mono-ampl:balance-capacidad}
Pour \(n\geq0\) et, dans la deuxième identité, \(n\geq1\),
\[
 \mathcal C_n=n-Z_n,\qquad
 \mathcal C_n-\mathcal C_{n-1}=1-z_n(0).
\]
\end{proposition}
\begin{proof}
L'irrationalité de \(\delta\) implique
\(\lceil(n+1)\delta\rceil=\lfloor(n+1)\delta\rfloor+1\).
Alors
\[
 \lfloor(n+1)(1-\delta)\rfloor
 =n+1-\lceil(n+1)\delta\rceil=n-Z_n.
\]
La soustraction de deux identités consécutives achève la preuve.
\end{proof}

Un événement \(z_n=1\) ajoute un bloc à l'histoire en maintenant la capacité
\(\mathcal C_n\) ; un événement \(z_n=0\) incrémente les deux compteurs. Sur
tout intervalle, les incréments de capacité et les lacunes ont pour somme
exactement sa longueur. En réduisant modulo \(q\), pour \(q=9\) ou
\(q=108\), on obtient
\[
 [\mathcal C_n]_q=[n]_q-[Z_n]_q.
\]
La phase uniforme des blocs et la phase de capacité sont liées par
la mémoire intégrée des lacunes. En particulier, une rétention de
capacité et un tour complet de la connexion ont des mécanismes
distincts, bien que tous deux puissent conserver une observation de phase.

L'inversion monotone de l'horloge donne un temps de première arrivée :
\[
 T_{\rm cap}(a)=\min\{n:\mathcal C_n=a\}
 =\left\lceil\frac a\rho\right\rceil-1,
 \qquad a\geq1.
\]
Avec \(\sigma=\rho^{-1}-1\), le temps nécessaire pour gagner \(b\) unités
de capacité est
\[
 T_{\rm cap}(a+b)-T_{\rm cap}(a)
 =b+\lceil(a+b)\sigma\rceil-\lceil a\sigma\rceil.
\]
La différence des plafonds prend l'un des deux entiers adjacents à
\(b\sigma\). On obtient ainsi \(9\) ou \(10\) blocs pour gagner
\(9\) unités de capacité, et \(113\) ou \(114\) pour en gagner \(108\).
Ce sont des temps de l'horloge inverse, mesurés depuis une première arrivée ; le
compteur des tours s'interprète toujours dans l'indice auquel il appartient.

\subsection{Suspension symbolique et cristal temporel apériodique}
\label{x_reciprocidad:mono-ampl:cristal}

Le codage par rotation réunit périodicité de phase, récurrence locale
et apériodicité globale. Soit \(\mathcal X_\delta\subseteq\{0,1\}^{\mathbb Z}\)
l'adhérence, dans la topologie produit, des translations du mot
bilatéral \(z_n(\theta)\). Notons \(\sigma_{\rm sh}\) le
décalage de ses suites. En conservant une phase nonadique, le pas est
\[
 F(r,z)=([r+1]_9,\sigma_{\rm sh}z),\qquad
 (r,z)\in C_9\times\mathcal X_\delta.
\]

\begin{definition}[Cristal temporel apériodique symbolique]
\label{x_reciprocidad:mono-ampl:def-cristal}
La suspension à toit unitaire du système précédent est l'espace
\[
 \mathcal S_\delta=
 \bigl((C_9\times\mathcal X_\delta)\times\mathbb R\bigr)/\sim,
 \qquad (y,t+1)\sim(Fy,t),
\]
avec le flot \(\Phi_s[y,t]=[y,t+s]\). On appelle \emph{cristal temporel
apériodique symbolique} ce modèle d'ordre temporel, accompagné de son
codage de lacunes et de la représentation du compteur des tours.
\end{definition}

L'adjectif temporel désigne l'action du flot et de sa section de retour ;
l'apériodicité caractérise ses itinéraires. La fréquence des uns dans chaque
suite est \(\delta\) : la borne uniforme de discrépance inférieure à un,
obtenue par télescopage dans toute fenêtre, persiste lors du passage à l'adhérence.
Une suite périodique aurait une fréquence rationnelle. Par conséquent,
\(\mathcal X_\delta\) ne contient pas d'orbites périodiques. La suspension
ne possède pas non plus d'orbites fermées : un retour du flot exigerait un temps
entier et une périodicité de l'itinéraire.

Les mots finis, en revanche, réapparaissent. Les itinéraires de longueur \(m\)
proviennent d'une partition de la circonférence par \(m+1\) points ; chaque
intervalle est visité de manière récurrente par la rotation irrationnelle. Même
lorsque les temps sont restreints à une phase fixée, le pas \(9\delta\) demeure
irrationnel. C'est pourquoi chaque mot apparaît dans les
\(9\) phases, et pourquoi valent les complexités déjà établies
\[
 p(m)=m+1,\qquad p_9(m)=9(m+1),\qquad p_3(m)=3(m+1).
\]
La récurrence des configurations locales coexiste avec l'absence d'une
période globale. Leur entropie topologique est nulle parce que ces complexités
croissent linéairement.

La réalisation par les observables \(V_9,W_\delta\) et le décalage
\(T\) exprime le même ordre par les covariances
\eqref{x_reciprocidad:eq:weyl-relojes}. Les fréquences de l'horloge uniforme appartiennent à
\(\frac19\mathbb Z+\delta\mathbb Z\), modulo \(1\). L'horloge de
capacité conserve également la relation
\([\mathcal C_n]_9=[n]_9-[Z_n]_9\) ; elle n'est pas remplacée par la phase uniforme
lors de l'interprétation des retours. Le terme \emph{cristal temporel} est ainsi
défini au moyen d'un système dynamique mathématique. Sa réalisation physique
requiert la spécification d'une dynamique énergétique, de ses observables temporels
et de sa stabilité, données distinctes de la présente construction symbolique.

\subsection{Symétrie spéculaire des régions et double réalisation}
\label{x_reciprocidad:mono-ampl:regiones}

L'échange spéculaire agit sur le bloc régional
\(B=(u^\pi,u^e,u^\varphi)\in(\mathbb F_3^6)^3\) par
\[
 \mathfrak cB=(u^e,u^\pi,u^\varphi).
\]
Son axe fixe et son agrégat sont respectivement \(u^\varphi\) et
\(s=u^\pi+u^e\). Cette fixation concerne l'involution dans une fibre ;
les deux coordonnées peuvent évoluer pendant le transport nonadique.
Avec \(q=u^\pi+u^e+u^\varphi\),
\(a=u^\pi+u^e-u^\varphi\) et \(d=u^\pi-u^e\), on retrouve
\[
 u^\varphi=2(q-a),\qquad s=q-u^\varphi,\qquad
 u^\pi=2(s+d),\qquad u^e=2(s-d).
\]
Les égalités sont composante par composante dans \(\mathbb F_3\), où
\(2^{-1}=2\). Le miroir conserve \(q,a,s\) et inverse \(d\) ; cette
coordonnée antisymétrique individualise la répartition que l'agrégat conserve
sans la distinguer. L'inversion \(\mathcal I\) de l'indice temporel et
l'involution \(\mathfrak c\) des canaux agissent donc sur des
données différentes.

Le retour de phase possède un témoin explicite :
\[
 \begin{aligned}
 B_1&=(012222\mid121221\mid112202),\\
 B_{10}&=(101222\mid112221\mid112002).
 \end{aligned}
\]
Les deux positions ont la phase \(1\) ; le bloc, l'axe et la mémoire ont
changé. De même, la synchronisation des dénombrements aux phases \(8\) et
\(9\) annule leurs différences relatives, tandis que la composante diagonale
passe de \((59,59,59)\) à \((43,43,43)\). L'égalité d'un observable
relatif n'épuise pas l'état qui le détermine.

La double réalisation conserve cette architecture. Dans la section
\ref{x_reciprocidad:exc:doble-lectura}, le même préfixe détermine d'une part des cylindres
emboîtés et d'autre part une suite de mots codés avec leurs repères.
La première application passe à la limite par contraction des diamètres ; la seconde
satisfait \(r_{n,m}Q_n=Q_mp_{n,m}\) et définit \(Q_\infty\).
La figure \ref{x_reciprocidad:mono-ampl:fig-regional} anticipe ces deux applications.
Le parcours incidenciel ultérieur construit des codes, des réseaux et la
réalisation de Moonshine dans leurs domaines propres ; le groupe des automorphismes
agit sur cette réalisation, et non sur les chiffres par une identification
implicite.


% BEGIN OWNER /Users/ruben/Documents/New project/output/EXTREMA_MEDIA_RAZON_RECIPROCIDAD_ES_EN_20260918/ARTICULO/ES/source/figures/monodromia_regional_ampliacion.tex
% Figura acumulativa; destino figures/monodromia_regional_ampliacion.tex.
\begin{figure}[p]
\centering
\begin{tikzpicture}[x=1cm,y=1cm,>=Latex,font=\small,
  flow/.style={->,line width=.55pt},
  ann/.style={align=center,font=\scriptsize}]
\node[font=\bfseries] at (6.9,9.2)
  {Connexion nonadique et coordonnées régionales};
\begin{scope}[shift={(3.4,6.8)}]
 \foreach \k in {1,...,9}{
  \pgfmathsetmacro{\ang}{90-40*(\k-1)}
  \coordinate (p\k) at (\ang:1.52);
  \fill (p\k) circle (1.4pt);
  \node[fill=white,inner sep=1.1pt,font=\scriptsize]
    at (\ang:1.83) {\(\k\)};
 }
 \foreach \k in {1,...,8}{
  \pgfmathtruncatemacro{\next}{\k+1}
  \draw[flow,shorten >=3pt,shorten <=3pt] (p\k)--(p\next);
 }
 \draw[flow,shorten >=3pt,shorten <=3pt] (p9)--(p1);
 \node[ann] at (0,.1) {\(g_k\in C_9\)\\\(\Gamma_{9,k}\)};
 \node[ann] at (0,-2.23) {Un tour : \(g\mapsto g\), \(m\mapsto m+1\)};
\end{scope}
\node[ann,text width=5.25cm] at (10.1,7.75)
 {\(B_k=(u_k^\pi,u_k^e,u_k^\varphi)\)\\
 bloc régional dans une fibre};
\node at (8.55,6.5) {\(u^\pi\)};
\node at (11.65,6.5) {\(u^e\)};
\draw[<->,line width=.65pt] (8.95,6.5)--node[above]{\(\mathfrak c\)}(11.25,6.5);
\draw[densely dashed,gray] (10.1,5.72)--(10.1,7.23);
\node[fill=white,inner sep=2pt] at (10.1,5.6){\(u^\varphi\)};
\node[ann] at (10.1,4.6)
 {Axe \(u^\varphi\) et agrégat \(u^\pi+u^e\)\\
 invariants sous \(\mathfrak c\)};
\draw[flow] (5.45,6.8)--(7.25,6.8);
\draw[gray!60,line width=.3pt] (.55,4.02)--(13.25,4.02);
\node[ann,text width=12cm] (hist) at (6.9,3.47)
 {Histoire compatible : préfixes, cylindres, orientation, quotients,
 repères d'incidence et mémoire};
\coordinate (fork) at (6.9,2.85);
\draw[flow] (hist.south)--(fork);
\draw[flow] (fork)--(3.6,2.2);
\draw[flow] (fork)--(10.3,2.2);
\node[ann,text width=5.6cm] at (3.6,1.65)
 {Réalisation archimédienne\\
 \(I_{n+1}\subseteq I_n,\quad |I_n|=729^{-n}\)\\
 limite des préfixes régionaux};
\node[ann,text width=5.6cm] at (10.3,1.65)
 {Réalisation incidencielle\\
 \(r_{n,m}Q_n=Q_mp_{n,m}\)\\
 mots codés et repères compatibles};
\node[ann,text width=5.6cm] at (3.6,.42)
 {\(\pi,\varphi,e\) ; détermination dodécaphasique de \(\alpha\)};
\node[ann,text width=5.6cm] at (10.3,.42)
 {Witt--Golay--Leech--\(V^\natural\)\\
 automorphismes de la réalisation exceptionnelle};
\draw[flow] (3.6,1.03)--(3.6,.78);
\draw[flow] (10.3,1.03)--(10.3,.78);
\end{tikzpicture}
\caption{Neuf phases d'une connexion et deux réalisations de l'histoire
compatible. Le cycle représente la phase ; chaque tour incrémente le compteur.
Le diagramme régional représente l'involution \(\mathfrak c\), qui
échange clôture et propagation et fixe l'axe d'autoéchelle dans une fibre.
Cette invariance n'impose pas que l'axe reste constant pendant le transport.
Les deux applications inférieures utilisent le même préfixe : l'une détermine
sa limite numérique et l'autre conserve les mots et les repères d'incidence. Les
applications de cette seconde réalisation sont spécifiées à la section
\ref{x_reciprocidad:exc:doble-lectura} ; le parcours exceptionnel est développé ensuite.}
\label{x_reciprocidad:mono-ampl:fig-regional}
\end{figure}

% END OWNER



% BEGIN OWNER /Users/ruben/Documents/New project/output/EXTREMA_MEDIA_RAZON_RECIPROCIDAD_ES_EN_20260918/ARTICULO/ES/source/sections/recta_nonadica_recuperada.tex
\clearpage
\subsubsection{Représentation rectiligne des phases et du retour}
\label{x_reciprocidad:mono-recta:desarrollo}

La représentation sur le segment \([0,10]\) conserve l'origine positionnelle d'APP et ordonne ses neuf marques intérieures. Dans ce diagramme, les entiers \(1,\ldots,9\) sont des représentants de phase. La position horizontale indique l'ordre des cartes ; les étiquettes supérieures spécifient des opérations régionales. La figure \ref{x_reciprocidad:mono-recta:figura} réunit ainsi la droite de départ, la fixation de l'axe d'auto-échelle, la distribution spéculaire des deux canaux et le retour avec mémoire (section~\ref{x_reciprocidad:mono-ampl:conexion}).


% BEGIN OWNER /Users/ruben/Documents/New project/output/EXTREMA_MEDIA_RAZON_RECIPROCIDAD_ES_EN_20260918/ARTICULO/ES/source/figures/recta_nonadica_0_10.tex
% Recuperación de fig_nonadic_route_mini.tex de la síntesis académica.
% Los rótulos designan operaciones regionales; la suma se efectúa en F_3^6.
\begin{figure}[H]
\centering
\begin{tikzpicture}[x=1.12cm,y=1cm,font=\small,>=Latex,
  guide/.style={line width=.45pt,black!55}]
 \draw[line width=.8pt] (0,0)--(10,0);
 \foreach \k in {0,...,10}{
  \draw[guide] (\k,-.09)--(\k,.09);
  \node[below=2.5mm] at (\k,0) {$\k$};
 }
 \foreach \k in {1,...,9}{\fill (\k,0) circle (1.7pt);}
 \node[anchor=west,font=\footnotesize] at (0,3.05)
   {Segment générateur $[0,10]$ et neuf représentants de phase};
 \draw[guide] (5,.12)--(5,1.7);
 \node[above,align=center,font=\footnotesize] at (5,1.7)
   {fixation de l'axe\\$u^\varphi$};
 \draw[decorate,decoration={brace,amplitude=4pt},guide]
   (6,.65)--(8,.65);
 \node[above=5pt,align=center,font=\footnotesize] at (7,.65)
   {agrégat régional\\$u^e+u^\pi$};
 \draw[guide] (9,.12)--(9,1.7);
 \node[above,align=center,font=\footnotesize] at (9,1.7)
   {involution spéculaire\\$\pi\leftrightarrow e$};
 \draw[->,line width=.7pt] (9,-.7)
   .. controls (9,-1.55) and (1,-1.55) .. (1,-.7);
 \node[fill=white,inner sep=3pt,font=\footnotesize] at (5,-1.23)
   {retour $9\to1$;\quad $m\mapsto m+1$};
 \node[font=\footnotesize] at (5,-1.95)
   {$w_6\to w_{12}\to w_{18}\to w_{24}\to w_{30}\to R_{36}\to G_9$};
\end{tikzpicture}
\caption{Représentation rectiligne de la connexion nonadique. La marque
\(5\) situe la première saturation forte et l'axe régional d'autoéchelle ;
l'accolade \(6\)--\(8\) identifie l'agrégat ternaire
fixé dans chaque fibre ; la phase \(9\) sélectionne l'orientation de
clôture. Le retour conserve la phase et incrémente la mémoire. Les
abscisses numérotent les positions opératives, tandis que les étiquettes
\(u^\pi,u^e,u^\varphi\) désignent des blocs de \(\mathbb F_3^6\).}
\label{x_reciprocidad:mono-recta:figura}
\end{figure}

% END OWNER


La fixation de l'axe se formule dans la fibre d'une frontière donnée. Si
\(q=u^\pi+u^e+u^\varphi\) et
\(a=u^\pi+u^e-u^\varphi\), alors
\[
 u^\varphi=2(q-a),\qquad s=u^\pi+u^e=q-u^\varphi
 \quad\text{dans }\mathbb F_3^6.
\]
Les données \((q,a)\) déterminent l'axe et l'agrégat ; la répartition ordonnée de \(s\) conserve la liberté spéculaire résolue par le prolongement. La cinquième phase possède une solution locale et une extension, avec la donnée de frontière \(c=111111\) et la signature résiduelle \(\rho_W=000\). C'est le sens de sa première saturation forte. Le symbole \(\varphi\) situé sur la marque \(5\) identifie l'axe régional qui intervient dans cette opération.

Les phases intermédiaires conservent cette décomposition au sein de chaque fibre,
tandis que leurs données de frontière évoluent. Le recensement complet des
multiplicités locales et des extensions de la branche distinguée est
\[
\begin{array}{c|rrrrrrrrr}
\text{phase}&1&2&3&4&5&6&7&8&9\\\hline
N_{\rm loc}&3&2&5&4&1&1&3&3&2\\
N_{\rm ext}&2&5&4&1&1&3&3&2&1
\end{array}
\]
La phase \(6\) combine unicité locale et trois prolongations ; la phase
\(7\) résout une fibre spéculaire triple ; la phase \(8\) synchronise
les trois recensements régionaux. La phase \(9\) conserve deux distributions
locales ayant le même axe et le même agrégat, dont l'une se
prolonge. Ces dénombrements expriment des fonctions différentes au sein d'une
même connexion et justifient les marques de la droite.

En particulier, pour les phases \(6,7,8\), les sommes régionales sont
respectivement \(012100\), \(101001\) et \(112000\). La conservation
représentée par l'accolade est une invariance par rapport à la répartition
spéculaire dans chaque fibre. L'évaluation réelle ultérieure applique les
lecteurs des cylindres compatibles aux histoires complètes ; la
somme interne représentée dans la figure appartient au domaine ternaire.

Enfin, l'arc \(9\to1\) représente le passage d'un tour au suivant. Les formules de la section \ref{x_reciprocidad:mono-ampl:conexion} conservent le reste de phase et incrémentent le compteur de tours. Le diagramme rectiligne et le diagramme cyclique de la figure \ref{x_reciprocidad:mono-ampl:fig-regional} expriment donc deux aspects complémentaires de la même holonomie : l'ordre des positions et la composition du retour sur l'état avec mémoire.

% END OWNER


\subsection{Transport temporel et composition exponentielle tritique}
\label{x_reciprocidad:mono-ampl:euler}

Les relations d'Euler de la section \ref{x_reciprocidad:euler-trit-articulacion}
décrivent les réalisations locales des trois régimes. Pour une signature
fixée \(\tau\), le générateur satisfait \(J_\tau^2=-\tau I\) ;
la conjugaison algébrique envoie \(J_\tau\) sur \(-J_\tau\). Par conséquent,
\[
 \exp(sJ_\tau)\exp(tJ_\tau)=\exp((s+t)J_\tau),\qquad
 \overline{\exp(tJ_\tau)}=\exp(-tJ_\tau),
\]
et le produit d'une évolution par sa conjuguée est l'unité. Dans le type
elliptique, cette réalisation est donnée par \(\cos^2t+\sin^2t=1\) ; dans l'hyperbolique,
par \(\cosh^2t-\sinh^2t=1\) ; dans le parabolique,
par \((I+tN)(I-tN)=I\), puisque \(N^2=0\). La même loi de conservation
admet trois expressions quadratiques.

La réalisation elliptique ferme un tour ; la nilpotente conserve un
déplacement linéaire ; l'hyperbolique sépare expansion et contraction avec
un déterminant unitaire. Aux paramètres régionaux déjà générés,
\[
 \exp(\pi J_{\mathrm e})=-I,\qquad
 \exp(J_{\mathrm p})=I+J_{\mathrm p},\qquad
 \exp(2\log\varphi\,J_{\mathrm h})=F_{\rm av}^{\,2}.
\]
L'évaluation \(\exp(I)=eI\) reconnaît l'unité de la loi exponentielle
commune. Le nombre \(e\) n'est pas attribué exclusivement au régime parabolique.

Le transport d'une réalisation locale par un isomorphisme \(U\) conserve
ces relations : si \(J'=UJU^{-1}\), la série exponentielle donne
\(U\exp(tJ)=\exp(tJ')U\). L'identité transporte le régime et sa
norme. Une transition entre types quadratiques différents appartient, en
revanche, à la sélection des feuillets et des régimes du TPK ; ce n'est pas une conjugaison
qui modifie la valeur de \(J^2\). Le produit temporel conserve donc
l'ordre des transitions et l'information de leurs extrémités.

Cette articulation distingue la clôture d'une représentation locale de
l'holonomie d'une histoire. On peut avoir \(\exp(2\pi J_{\mathrm e})=I\)
en même temps que \(S|r,m\rangle=|r,m+1\rangle\). La première égalité
rétablit l'orientation du régime ; la seconde enregistre un autre tour. La
structure discrète conserve les deux : le retour local et la provenance
accumulée de l'état qui le réalise.

% END OWNER



% BEGIN OWNER /Users/ruben/Documents/New project/output/EXTREMA_MEDIA_RAZON_RECIPROCIDAD_ES_EN_20260918/ARTICULO/ES/source/sections/primos_retornos.tex
% Recuperación del producto nativo y su selector, no incorporación de RH.
\subsection{Irréductibilité multiplicative et périodes de retour}
\label{x_reciprocidad:primos-retornos}

La distinction entre structure, évaluation et phase apparaît aussi dans le
secteur multiplicatif. Le produit local d’APP contient à la fois le reste
et le quotient ; son itération permet de construire des formes normales et
des factorisations avant d’examiner leurs périodes décimales. Ce développement
situe l’expression \emph{mode premier} dans un domaine précis et relie
l’irréductibilité aux opérations déjà définies, non à une coïncidence
isolée de périodes (section~\ref{x_reciprocidad:primos-retornos}).

Soit
\[
 \mathcal W_9=\{\varnothing\}\sqcup
 \bigsqcup_{\ell\geq1}\{1,\ldots,9\}^{\ell}.
\]
Les suites sont ordonnées à partir du chiffre le moins significatif. L'évaluation
\(\operatorname{ev}_9(u)=\sum_j u_j9^j\), avec
\(\operatorname{ev}_9(\varnothing)=0\), est bijective sur les entiers non
négatifs : pour un entier positif, la division de \(n-1\) par \(9\)
détermine de manière unique le premier chiffre et la continuation. C'est la
numération nonadique bijective ; elle conserve le représentant \(9\) de la
frontière.

La somme \(\boxplus_\#\) s'obtient en appliquant le feuillet additif d'APP aux chiffres de chaque position et en propageant leurs quotients. S'il n'y a qu'un chiffre, celui-ci constitue le résidu provisoire ; un quotient entrant non nul est normalisé au moyen d'une seconde opération additive. La somme des quotients est transportée à la position suivante. Pour multiplier une suite par un chiffre \(d\), la cellule multiplicative produit \(u_jd=9q_j+r_j\), et le quotient entrant \(c_j\) est incorporé au moyen de
\[
 (w_j,c_{j+1})=
 \begin{cases}
 (r_j,q_j),&c_j=0,\\
 \bigl(R^+(r_j,c_j),q_j+q^+(r_j,c_j)\bigr),&c_j>0.
 \end{cases}
\]
On commence avec \(c_0=0\) et on incorpore le quotient final lorsqu'il est non
nul. Pendant l'addition, \(c_j\leq2\) ; pendant la multiplication par un chiffre,
\(c_j\leq9\). Les opérations locales reçoivent ainsi des arguments dans
leur domaine. Notons \(d\odot_\#u\) cette seconde procédure,
avec \(d\odot_\#\varnothing=\varnothing\).

Pour \(v=(v_0,\ldots,v_{m-1})\), le produit complet est construit par
la récurrence
\[
 a_m=\varnothing,\qquad
 a_j=(9\odot_\#a_{j+1})\boxplus_\#(v_j\odot_\#u),
 \qquad u\boxtimes_\#v=a_0.
\]
La règle compose les deux feuillets locaux. Son évaluation vérifie
\begin{equation}
 \operatorname{ev}_9(u\boxtimes_\#v)
   =\operatorname{ev}_9(u)\operatorname{ev}_9(v).
 \label{x_reciprocidad:eq:primos-entrelazamiento}
\end{equation}
En effet, chaque transition conserve la somme partielle augmentée du quotient
en attente pondéré par la puissance de \(9\). La récurrence donne alors
\(\operatorname{ev}_9(a_j)=9\operatorname{ev}_9(a_{j+1})+
v_j\operatorname{ev}_9(u)\), dont l'itération prouve
\eqref{x_reciprocidad:eq:primos-entrelazamiento}. L'unicité de la forme normale transporte
l'associativité et la commutativité vers le produit construit ; son unité est \((1)\).

\begin{definition}[Coproduit multiplicatif réduit]
Sur l'espace vectoriel à support fini engendré par les formes normales non vides, on définit
\[
 \widetilde\Delta_\# e_w
 =\sum_{\substack{u\boxtimes_\#v=w\\u,v\ne(1)}}e_u\otimes e_v.
\]
Une forme normale non unitaire est irréductible lorsque son image par
\(\widetilde\Delta_\#\) est nulle.
\end{definition}

Les fibres de factorisation sont finies d'après
\eqref{x_reciprocidad:eq:primos-entrelazamiento}. Dans la complétion
\(\mathcal H_\#=\ell^2(\mathcal W_9\setminus\{\varnothing\})\),
les images de formes distinctes ont des supports disjoints. Si
\(d_\#(w)\) compte les factorisations ordonnées non unitaires, le domaine
de la clôture est
\[
 \operatorname{Dom}(\overline{\widetilde\Delta_\#})
 =\left\{x\in\mathcal H_\#:
       \sum_w d_\#(w)|x_w|^2<\infty\right\}.
\]
En effet, le carré de la norme de l'image est cette somme pondérée. Par conséquent, l'opérateur
\[
 A_\#=(\overline{\widetilde\Delta_\#})^*
                  \overline{\widetilde\Delta_\#}
\]
est diagonal, positif et autoadjoint : sa valeur propre en \(e_w\) est le
nombre de factorisations ordonnées non unitaires de \(w\). La projection
\[
 P_{\rm irr}=(I-|e_{(1)}\rangle\langle e_{(1)}|)
                    \mathbf1_{\{0\}}(A_\#)
\]
sélectionne exactement les formes irréductibles. Par l'évaluation multiplicative, celles-ci correspondent aux nombres premiers ordinaires. La sélection dépend de la structure de factorisation et maintient séparé l'état unitaire.

Pour une valeur \(n\geq2\) première avec \(10\), la division décimale induit la
permutation des résidus \(r\mapsto10r\pmod n\). Le retour du résidu
\(1\) a pour période
\[
 \tau_n=\operatorname{ord}_n(10).
\]
Si \(n=\prod_p p^{a_p}\), le théorème chinois des restes donne
\[
 \tau_n=\operatorname{mcm}_{p^{a_p}\parallel n}
                    \operatorname{ord}_{p^{a_p}}(10).
\]
Le retour composé synchronise les retours des puissances premières.
Pour un nombre premier \(p\notin\{2,3,5\}\), on a
\[
 \tau_p\mid p-1,\qquad
 M_p=\frac{10^{\tau_p}-1}{p}\in\mathbb N,
 \qquad M_p\equiv0\pmod9.
\]
La dernière congruence exprime la clôture nonadique de la répétition décimale : \(10^{\tau_p}-1\) est un multiple de \(9\) et \(p\) est inversible modulo \(9\). La même congruence vaut pour tout \(n\) premier avec \(90\). En particulier, \(7,13,77,91,143\) ont pour période décimale \(6\), bien que seuls les deux premiers soient irréductibles. Le coproduit et la période enregistrent donc des propriétés différentes d'une même valeur construite.

Cette différenciation précise la relation avec le raffinement apériodique.
La suite des vacances détermine les incréments de l'indice de
capacité ; la factorisation détermine les modes multiplicatifs ; la
division décimale détermine leurs temps de retour. Le représentant
nonadique conserve la fermeture de phase, tandis que le quotient et la
factorisation conservent une information que cette fermeture ne distingue pas. Le
lien avec les constructions spectrales du traité réside dans ces
opérateurs et leurs entrelacements. La démonstration de leurs résultats
analytiques ultérieurs conserve son développement propre ; elle n'est pas remplacée par
la périodicité d'une observation résiduelle.

\subsubsection{Discrépance des vacances et lectures de Dirichlet}

Le lien admet également une expression analytique exacte. On conserve le
même codage \(z_n=z_n(0)\), pour \(n\geq1\), et on considère
deux observables ultérieurs : la série sur tous les indices et le
produit sur les indices sélectionnés par irréductibilité. Pour
\(\operatorname{Re}s>1\),
\[
 Z_{\rm vac}^{+}(s)=\sum_{n\geq1}\frac{z_n}{n^s},\qquad
 Z_{\rm vac}^{\times}(s)=\prod_p(1-p^{-s})^{-z_p}.
\]
Tous deux convergent absolument dans ce demi-plan. Leur facteur commun est la
densité \(\delta\), mais leurs discrépances sont différentes.

\Needspace{21\baselineskip}
\begin{proposition}[Séparation de la densité de vacances]
On a
\[
 Z_{\rm vac}^{+}(s)=\delta\zeta(s)+H_{\rm vac}(s),
 \qquad H_{\rm vac}\in\mathcal O(\operatorname{Re}s>0).
\]
Dans \(\operatorname{Re}s>1\), avec les logarithmes définis par les séries absolument convergentes d'Euler,
\[
 Z_{\rm vac}^{\times}(s)=\zeta(s)^\delta G_{\rm vac}(s),\qquad
 \log G_{\rm vac}(s)=
 \sum_p(z_p-\delta)\log(1-p^{-s})^{-1}.
\]
\end{proposition}
\begin{proof}
Les sommes partielles centrées satisfont
\[
 A_N:=\sum_{n=1}^N(z_n-\delta)
 =\lfloor(N+1)\delta\rfloor-N\delta,\qquad |A_N|<1.
\]
La sommation d'Abel représente le terme centré par \(s\int_1^\infty A_{\lfloor x\rfloor}x^{-s-1}\,dx\). L'intégrale converge localement uniformément dans \(\operatorname{Re}s>0\), ce qui prouve l'holomorphie. La seconde identité résulte de la séparation de \(z_p=\delta+(z_p-\delta)\) dans le logarithme du produit, au sein de son domaine de convergence absolue.
\end{proof}

La décomposition rend visible ce qu'ajoute la lecture sur les nombres premiers : elle examine
la même suite de vacances sur un sous-ensemble déterminé par la
factorisation. Le contrôle de la discrépance sur tous les entiers donne
le premier prolongement holomorphe ; le second observable conserve une
discrépance propre sur les nombres premiers. Il s'agit de deux applications de la même
structure discrète, avec des domaines analytiques explicites. Cette articulation
explique leur pertinence dans le noyau arithmétique sans transférer dans cet
article la démonstration des résultats terminaux du traité.

% END OWNER

